\documentclass[11pt,a4paper]{article}

\usepackage{fontspec}
\setmainfont{Arial}[Path=/System/Library/Fonts/Supplemental/, Extension=.ttf, UprightFont=Arial, BoldFont=Arial Bold, ItalicFont=Arial Italic, BoldItalicFont=Arial Bold Italic]
\usepackage[bulgarian]{babel}
\usepackage[a4paper,margin=2.2cm]{geometry}
\usepackage{amsmath}
\usepackage{amssymb}
\usepackage{booktabs}
\usepackage{array}
\usepackage{tabularx}
\usepackage{xltabular}
\usepackage{makecell}
\usepackage{multirow}
\usepackage{colortbl}
\usepackage{enumitem}
\usepackage[table]{xcolor}
\usepackage{tikz}
\usetikzlibrary{arrows.meta,positioning,shapes.geometric,fit,calc,backgrounds}
\usepackage{forest}
\usepackage{fancyhdr}
\usepackage{titlesec}
\usepackage{hyperref}

\definecolor{navy}{HTML}{1F3A5F}
\definecolor{teal}{HTML}{2A7F7F}
\definecolor{sand}{HTML}{F3EFE6}
\definecolor{riskL}{HTML}{D8EFD3}
\definecolor{riskM}{HTML}{FFF1C2}
\definecolor{riskH}{HTML}{F8D0CC}
\definecolor{grayline}{HTML}{C9CED6}

\hypersetup{colorlinks=true, linkcolor=navy, urlcolor=navy, pdftitle={Курсов проект: STEM център в СМГ}, pdfauthor={Алекс Цветанов}}

\setlength{\parindent}{0pt}
\setlength{\parskip}{5pt}
\setlist{nosep, leftmargin=1.4em}
\renewcommand{\arraystretch}{1.15}

\titleformat{\section}{\Large\bfseries\color{navy}\raggedright\hyphenpenalty=10000}{\thesection.}{0.5em}{}
\titleformat{\subsection}{\large\bfseries\color{navy}\raggedright\hyphenpenalty=10000}{\thesubsection.}{0.5em}{}
\titlespacing*{\section}{0pt}{6pt}{6pt}
\titlespacing*{\subsection}{0pt}{8pt}{3pt}

\pagestyle{fancy}
\fancyhf{}
\fancyhead[L]{\small Курсов проект: интердисциплинарен STEM център в СМГ}
\fancyhead[R]{\small Управление на проекти, 2026/2027}
\fancyfoot[C]{\thepage}
\renewcommand{\headrulewidth}{0.4pt}

% Курсови теми от предишните лекции — малък етикет в текста
\newcommand{\lecturetag}[1]{\textcolor{teal}{\small\bfseries[#1]}}

% --- numbers.tex
% Генерирано от бюджета (gen_numbers.py). Всички суми са в € с ДДС.
\newcommand{\NumSTEM}{299\,998,58~€}
\newcommand{\NumFurn}{36\,039,17~€}
\newcommand{\NumEquip}{211\,934,26~€}
\newcommand{\NumSpec}{22\,258,45~€}
\newcommand{\NumSoft}{29\,766,71~€}
\newcommand{\NumVOSKS}{59\,984,12~€}
\newcommand{\NumSoftVOSKS}{1\,219,20~€}
\newcommand{\NumBoth}{359\,982,71~€}
\newcommand{\NumMgmt}{3\,599,83~€}
\newcommand{\NumTotal}{363\,582,54~€}
\newcommand{\PctUse}{99,995\,\%}
\newcommand{\NumItems}{103}
\newcommand{\NumLines}{133}
\newcommand{\NumBGItems}{80}
\newcommand{\NumOffers}{23}
\newcommand{\NumSellers}{53}
\newcommand{\NumEquipRob}{73\,806,22~€}
\newcommand{\NumEquipSci}{74\,495,72~€}
\newcommand{\NumEquipMath}{63\,632,32~€}
\newcommand{\NumFurnRob}{9\,097,04~€}
\newcommand{\NumFurnSci}{22\,591,31~€}
\newcommand{\NumFurnMath}{4\,350,82~€}
\newcommand{\NumSoftRob}{8\,625,24~€}
\newcommand{\NumSoftSci}{515,90~€}
\newcommand{\NumSoftMath}{20\,625,56~€}
\newcommand{\NumVOSKSEq}{58\,764,92~€}
\newcommand{\NumLeftSTEM}{1,42~€}
\newcommand{\NumLeftVOSKS}{15,88~€}
\newcommand{\NumLeftTotal}{17,46~€}

% --- numbers_extra.tex
\newcommand{\NumFlexBuffer}{22\,666,57~€}
\newcommand{\NumOffersEU}{18}
\newcommand{\NumOffersNonEU}{5}
\newcommand{\NumOfferTotal}{97\,161,20~€}
\newcommand{\PctOfferTotal}{26,9\,\%}
\newcommand{\TopLines}{преносим компютър за ученици (26~бр., 38\,390,35~€); интерактивен дисплей 75" 4K (8~бр., 23\,467,39~€); преносим компютър за ученици (13~бр., 19\,195,18~€); безжична цифрова лаборатория (14~бр., 16\,757,50~€); лабораторна ученическа маса (6~бр., 12\,184,70~€)}

\usepackage{caption}

\begin{document}

% --- sections/title.tex
\thispagestyle{empty}
\begin{center}
  {\small Технически университет – София \quad·\quad Факултет „Компютърни системи и технологии“\par}
  \vspace{2.2cm}
  {\Large\bfseries\color{navy} Курсов проект\par}
  \vspace{6pt}
  {\LARGE\bfseries Разработване и управление на интердисциплинарен\\[2pt] STEM център с интеграция на изкуствен интелект\par}
  \vspace{14pt}
  {\Large „СМГ STEM Хъб“ – Софийска математическа гимназия „Паисий Хилендарски“\par}
  \vspace{1.6cm}
  \begin{tabular}{@{}lp{0.66\textwidth}@{}}
    Учебна дисциплина: & Управление на проекти, 2026/2027 г. \\
    Студент:           & Алекс Иванов Цветанов, фак. № 121326006 \\
    Вид:               & индивидуален курсов проект (3 технологични направления) \\
    Училище:           & СМГ „Паисий Хилендарски“, ул. „Искър“ 61, София \\
    Направления:       & Математика и информатика · Роботика и кибер-физични системи · Природни науки \\
    Бюджет (с ДДС):    & \NumTotal{} от лимит 363\,600~€ (\PctUse) \\
  \end{tabular}
\end{center}

\vspace{1.2cm}
\noindent\colorbox{sand}{\parbox{\dimexpr\textwidth-2\fboxsep}{%
\textbf{Основен въпрос.} \emph{„Как чрез оптимална комбинация от технологии, учебни дисциплини, данни, физическа среда и изкуствен интелект може да бъде създаден ефективен интердисциплинарен STEM център?“}

\smallskip
\textbf{Кратък отговор.} Като трите лаборатории (данни и AI, роботика и мейкърспейс, природни науки) и осемте високотехнологични класни стаи не се оборудват поотделно, а се свързват в една среда около общ център за данни – „СМГ STEM Data Hub“. Всеки сензор, робот и компютър подава данни, които се използват едновременно в няколко предмета; три AI приложения (обучение, анализ на данни, управление на средата) работят върху тези данни локално в училището, под контрола на учител. Така едно и също оборудване обслужва седем учебни предмета и всички ученици от V до XII клас.

\smallskip
\textbf{Защо това решение е по-добро от алтернативите.} \emph{По-интегрирано} – общите данни и общият график свързват предметите, вместо всеки кабинет да работи отделно; \emph{по-ефективно} – споделеното оборудване се използва в повече часове и не се купува два пъти; \emph{по-устойчиво} – AI приложенията за учениците работят локално с отворени модели, без абонаменти, и личните данни на учениците не напускат училището; платените годишни лицензи са малко и са само за учителите. Многокритериалната оценка в раздел 8 дава на този вариант 4,56 от 5 точки срещу 2,07–3,60 за алтернативите.}}

\vfill
\begin{center}\small София, ноември 2026 г.\end{center}
\clearpage


\clearpage
{\hypersetup{linkcolor=black}\small\tableofcontents}
\clearpage% --- sections/s1.tex
\section{Дефиниране на проекта и SMART цели}

\textbf{Основен проблем.} СМГ има силен профил по математика, информатика и природни науки и олимпийска традиция, но STEM обучението е разпокъсано: предметите работят в отделни кабинети, уредите се ползват от отделни учители, роботиката е само извънкласен клуб, а данни от реални експерименти почти не се събират и не се споделят. AI инструментите се ползват стихийно в облака, без контрол върху данните на непълнолетните.

\subsection{Описание на идеята за изграждане на STEM среда}

\textbf{Направления и пространства.} Центърът обхваща три технологични направления в съществуващи помещения (без СМР~– само обзавеждане и оборудване): \textbf{Математика и информатика}~– лаборатория „Данни и AI“ (26 места); \textbf{Роботика и кибер-физични системи}~– лаборатория „Роботика и мейкърспейс“ в сутерена, където сега работи Клуб „Роботика“ (13 станции $\times$ 2 = 26 места и мейкър зона); \textbf{Природни науки}~– едноименна лаборатория (6 острова $\times$ 4 = 24 места) с умна оранжерия. Към тях се добавят \textbf{8 ВОСКС} в основната сграда (интерактивен дисплей с OPS, камера, документ-камера, лаптоп). \textbf{Интегриращата идея} е общият „СМГ STEM Data Hub“ (ThingsBoard с MQTT, NAS, локален AI сървър, Grafana, Moodle, резервации) вместо изолирани лаборатории: всеки сензорен поток се записва в едно хранилище, използва се от поне три предмета и е видим и във ВОСКС, а лаптопите, 3D принтерите и AI сървърът се споделят от трите направления по общ график.

\textbf{Допълнително специализирано оборудване се предвижда} (\NumSpec, до 20\,\% от оборудването), защото го налагат профилът на СМГ и опитът на учителите, които подготвят олимпийци и водят Клуб „Роботика“: AI работна станция (локален езиков модел~– данните на учениците остават в училището, GDPR), лазерен резач и 3D скенер (прототипи и шасита), спектрофотометър за видимата област, 335–800~nm (количествен анализ за олимпиадите по химия и биология), GoTo телескоп с астрокамера (астрономия).

\textbf{Учебни предмети (7).} Ядрото са \emph{математика, информатика, физика и астрономия, биология и здравно образование}~– профилни предмети на СМГ с олимпийска традиция; добавят се \emph{химия и опазване на околната среда, информационни технологии, география и икономика}, които използват същите данни (химичен анализ, табла, климат и картиране). Предметите се включват чрез 4–6-седмични сценарии, планирани съвместно от учителите: всеки предмет работи в своите часове (ЗУП, ИУЧ, клубове), но върху общите данни, а продуктът се защитава пред учителите по включените предмети. \emph{Пример~– „Умна оранжерия“:} по биология учениците поставят опит с растения при различно осветление; по физика калибрират сензорите; по химия следят pH и EC; по информатика програмират ESP32 и обмена по MQTT; по математика строят регресионен модел на растежа и го сравняват с прогнозата на AI модела (поливането е автоматично в безопасни граници, торенето~– само с потвърждение от учител).

\subsection{SMART цели}

\textbf{S (конкретна)} \lecturetag{Лекция 2: SMART}. До 15.09.2027~г. в СМГ да работи STEM център с три лаборатории, 8 ВОСКС, общ Data Hub и три AI приложения (тютор, анализ на данни от оранжерията, график и прогнозна поддръжка). Той заменя кабинетното обучение с проектно върху реални данни, облачния AI~– с локален, а дублираните уреди~– със споделени.

\textbf{M (измерима)}~– показатели за първата учебна година:

{\footnotesize
\begin{xltabular}{\textwidth}{@{}>{\raggedright\arraybackslash}p{4.3cm}>{\raggedright\arraybackslash}p{0.9cm}>{\raggedright\arraybackslash}X>{\raggedright\arraybackslash}X@{}}
\toprule
\textbf{Показател} & \textbf{База} & \textbf{Цел} & \textbf{Измерване / източник} \\
\midrule
Интегрирани предмети & 0 & $\geq 7$ & график на центъра, Moodle \\
Общи STEM сценарии & 0 & $\geq 8$ & сценарии в Moodle \\
Използване на оборудването & – & $\geq 70\,\%$ от слотовете & резервации + телеметрия \\
Автоматизирани процеси & 0 & $\geq 5$ (поливане, автоматично изключване и др.) & правила в ThingsBoard \\
AI приложения в употреба & 0 & 3; тютор с $\geq 90\,\%$ точност & логове, външен тест \\
Ученици едновременно & – & 76 в лабораториите + 8 паралелки & работни места, разписание \\
Икономия на ресурси & – & $-15\,\%$ енергия, без дубликати & консумация, инвентар \\
Резултати от STEM проекти & – & $\geq 30$, от тях $\geq 10$ на състезания & регистър в Moodle \\
\bottomrule
\end{xltabular}}

\textit{Методология} („–“~– няма данни): за всеки показател са определени формула, източник, отговорник и честота (напр. използване = заети $\div$ налични слотове; икономия~– спрямо октомври 2027~г.). Данните идват от Data Hub и резервациите, допълват се с анкета и се отчитат всеки срок пред Педагогическия съвет; отклонение над 10\,\% води до коригиращи мерки.

\textbf{A (постижима, разпределена, приета).} \emph{Финансово и технологично:} \NumSTEM{} за STEM (лимит 300\,000~€) и \NumVOSKS{} за ВОСКС (лимит 60\,000~€) по реални оферти за серийни продукти; в приоритетната матрица разходите се „ограничават“, обхватът се „подобрява“, а времето се „приема“ \lecturetag{Лекция 3: приоритетна матрица}. \emph{AI интеграцията} се приема само след \textbf{независимо външно тестване} (точност, сигурност, натоварване, fail-safe): 3\,300~€ от средствата за управление (\NumMgmt), 3 седмици през август 2027~г. и 4-седмичен пилот с 2 паралелки (15.09–12.10.2027~г.; т.~3.3). \emph{Хора} (т.~6): ръководител на проекта е зам.-директорът по учебната дейност; ключови са учителите по физика и информатика, ръководителят на Клуб „Роботика“ и системният администратор. \emph{Заинтересовани страни:} директорът подписва с КЕП; Педагогическият съвет одобрява; Националният STEM център (НСЦ) съгласува, МОН финансира; Общественият съвет~– становище; родителите~– съгласие за данните; учениците~– обратна връзка.

\textbf{R (релевантна).} Центърът развива изследователски умения, работа с данни, инженерен дизайн и AI грамотност чрез задачи, които не принадлежат на един предмет, при това в задължителните и избираемите часове~– за всички около 1300 ученици, а не само за олимпийците.

\textbf{T (обвързана със срок)} \lecturetag{Лекция 1: жизнен цикъл}. Центърът се проектира, оборудва и интегрира за 10 месеца, а проектът трае 21 месеца: иницииране~– ноември–декември 2026~г.; планиране~– януари–февруари 2027~г.; изпълнение~– ЗОП (март–май), доставка и монтаж (юни–юли), интеграция и външен тест на AI (август); експлоатация от \textbf{15.09.2027~г.}; мониторинг и контрол~– до юни 2028~г.; закриване~– юли 2028~г.

\clearpage% --- sections/s2.tex
\section{Модел за интердисциплинарна STEM интеграция}
% Отметка за таблицата (шрифтът Arial няма знак ✓)
\providecommand{\sIIcheck}{\tikz[baseline=-0.25ex,x=1ex,y=1ex]\draw[teal,line width=1.1pt,line cap=round,line join=round](0,0.5)--(0.4,0.05)--(1.15,1.05);}

Изискването STEM центърът да не бъде сбор от отделни лаборатории е заложено още при определянето на обхвата \lecturetag{Лекция 3: обхват и WBS}. Интеграцията е самостоятелен работен пакет в WBS („Интеграция на Data Hub, IoT и AI“), а не страничен резултат от доставките, и служи като критерий за обхвата: оборудване се включва в бюджета само ако (1)~се използва в поне два учебни предмета и (2)~е свързано с общия „СМГ STEM Data Hub“ – подава или визуализира данни или поне се резервира и отчита през общата система. Пространствата на модела са Л1 „Данни и AI“, Л2 „Роботика и мейкърспейс“, Л3 „Природни науки“ с умната оранжерия и 8~ВОСКС.

\subsection{STEM Integration Matrix}

Всеки ред на матрицата е верига от седем звена: \emph{учебна дисциплина → STEM област → технология → оборудване → данни → AI → учебна дейност}. Разработени са 14~вериги, групирани по направления (пълна таблица – Приложение~А; графична схема на връзките между дисциплини, технологии и пространства – Приложение~Б). Три от тях:
\begin{itemize}
\item \textbf{Биология} → природни науки → IoT сензори → умна оранжерия (Л3) → температура, влажност на почвата, светлина → AI-2 (прогноза на времеви редове и компютърно зрение) → прогнозиране на условията за развитие на растенията.
\item \textbf{Информатика} → програмиране → микроконтролери → ESP32/Arduino и Raspberry Pi~5 с AI HAT+ (Л2) → сензорни данни по MQTT → ML модел на самото устройство (edge AI) → автоматично управление на поливането в безопасни граници.
\item \textbf{Математика} → анализ на данни → локална вместо облачна платформа (ThingsBoard + NAS в училището) → исторически данни от оранжерията и метеостанцията → AI-2 → статистически и предиктивен анализ с оценка на грешката на прогнозата.
\end{itemize}
Матрицата покрива всичките 7~предмета (всеки участва в поне 6~вериги), трите направления и ВОСКС, трите AI приложения (AI-1 – AI-3) и осемте STEM сценария за първата година; тя и схемата са обобщени и на слайда за т.~2 в презентацията.

\subsection{Една технология – много учебни предмети}

Всяка технология обслужва поне три предмета, средно 5,1 от 7 (табл.~\ref{tab:s2reuse}), затова се купува веднъж за центъра, а не поотделно за всеки кабинет.

\begin{table}[h]
\centering\footnotesize
\setlength{\tabcolsep}{4pt}
\renewcommand{\arraystretch}{1.08}
\caption{Технология × учебен предмет (М – математика, И – информатика, ИТ – информ. технологии, Ф – физика и астрономия, Х – химия, Б – биология, Г – география)}
\label{tab:s2reuse}
\begin{tabularx}{\textwidth}{@{}>{\raggedright\arraybackslash}X l *{7}{c} c@{}}
\toprule
\textbf{Технология} & \textbf{Пространство} & \textbf{М} & \textbf{И} & \textbf{ИТ} & \textbf{Ф} & \textbf{Х} & \textbf{Б} & \textbf{Г} & \textbf{Σ} \\
\midrule
IoT сензорен поток (оранжерия, метеостанция, CO$_2$ сензори) & Л3, ВОСКС & \sIIcheck & \sIIcheck & \sIIcheck & \sIIcheck & \sIIcheck & \sIIcheck & \sIIcheck & 7 \\
AI сървър (локален LLM + компютърно зрение) & Data Hub & \sIIcheck & \sIIcheck & \sIIcheck & \sIIcheck & \sIIcheck & \sIIcheck & \sIIcheck & 7 \\
Интерактивни дисплеи с Grafana табла & 8 ВОСКС & \sIIcheck & \sIIcheck & \sIIcheck & \sIIcheck & \sIIcheck & \sIIcheck & \sIIcheck & 7 \\
Роботи, роботизирана ръка и дронове & Л2 & \sIIcheck & \sIIcheck & & \sIIcheck & \sIIcheck & \sIIcheck & \sIIcheck & 6 \\
3D принтери, лазерен резач, 3D скенер & Л2 & \sIIcheck & & \sIIcheck & \sIIcheck & & \sIIcheck & \sIIcheck & 5 \\
Логери, pH/EC метри, спектрофотометър & Л3 & \sIIcheck & & & \sIIcheck & \sIIcheck & \sIIcheck & & 4 \\
Микроконтролери и edge AI (ESP32, Raspberry Pi~5) & Л2 → Л3 & & \sIIcheck & \sIIcheck & \sIIcheck & & \sIIcheck & & 4 \\
Цифрови микроскопи с камера & Л3 & & \sIIcheck & & & \sIIcheck & \sIIcheck & & 3 \\
GoTo телескоп с астрокамера & Л3 & \sIIcheck & \sIIcheck & & \sIIcheck & & & & 3 \\
\bottomrule
\end{tabularx}
\end{table}

\textbf{Сензорният поток от оранжерията и метеостанцията} е най-широко споделеният ресурс: един и същ запис се използва в биологията (условия за растеж), химията (pH и EC на хранителния разтвор), физиката (топлообмен, светлинен поток), математиката (статистика, регресия, прогноза), информатиката (MQTT, бази от времеви редове, ML) и географията (микроклимат спрямо външния климат). \textbf{Роботите и дроновете} от Л2 се ползват във физиката (кинематика, сили, енергия на батериите), математиката (траектории, вектори, координати) и географията (картиране и площи на училищния двор), а роботизираната ръка пренася автоматизацията в химията („Робот-лаборант“). \textbf{AI сървърът} е общ за информатиката (обучение и оценка на модели), биологията (класификация на клетки), химията (разпознаване на точката на еквивалентност) и чрез AI тютора – за всички седем предмета. Сървърът е един и локален – данните на учениците не напускат училището (GDPR).

\subsection{Механизми срещу изолираното функциониране на лабораториите}

\begin{itemize}
\item \textbf{Общ Data Hub.} Всички устройства – от логерите в Л3 до CO$_2$ сензорите във ВОСКС – публикуват по MQTT в ThingsBoard по единна схема (пространство, устройство, величина, време); историята в NAS е достъпна за всеки предмет, не само за кабинета-източник.
\item \textbf{Обща система за резервации и инвентар.} Оборудването е ресурс на центъра, а не на кабинет: мобилните комплекти (логери, роботи, дронове) се резервират за Л1–Л3 и ВОСКС по общ график, който AI-3 предлага, а администраторът одобрява; заетите слотове се отчитат като степен на използване (цел – поне 70\,\%).
\item \textbf{ВОСКС като прозорец към лабораториите.} Дисплеите показват на живо Grafana таблата и видеото от микроскопа и телескопа – паралелка в обикновена класна стая работи с реални данни, а самите ВОСКС са източник на данни („Климатът на училището“).
\end{itemize}
Интеграцията се измерва със SMART показателите от т.~1 \lecturetag{Лекция 2: SMART} (предмети, сценарии, използване на оборудването) и с критерии 1–5 на STEM Integration Score (т.~8).

\textbf{Външно тестване на AI интеграцията.} Всяка от 14-те вериги минава през AI звено, затова връзката сензор → AI → действие не се приема „на доверие“, а след независимо външно приемно тестване (цена, срок и критерии – т.~3.3).

\clearpage% --- sections/s3.tex
\section{AI интеграция в STEM центъра}

Предлагат се три AI приложения – по едно за всяка от трите области, посочени в заданието (обучение, анализ на данни, управление на средата). Всяко има ясно място в центъра: AI тюторът се ползва най-вече в Лаборатория „Данни и AI“, анализът на данни – в Лаборатория „Природни науки“ и умната оранжерия, а управлението на средата обхваща трите лаборатории и 8-те ВОСКС. И трите работят върху „СМГ STEM Data Hub“ (AI работна станция, NAS, ThingsBoard, Grafana) с отворени модели, изпълнявани в училището, така че личните данни на учениците не напускат сградата (GDPR).

\subsection{Предложени AI приложения и схема на функциониране}

\textbf{AI-1 „СМГ AI тютор“} (AI за обучение) помага при програмиране, генерира задачи според нивото на ученика и дава обратна връзка по лабораторните протоколи. \textbf{AI-2 „Умна оранжерия и лабораторни данни“} (AI за анализ на данни) прогнозира микроклимата, открива аномалии и разпознава състоянието на растенията. \textbf{AI-3 „Интелигентен график и прогнозна поддръжка“} (AI за управление на средата) предлага график на лабораториите и ВОСКС, предупреждава за износване и посочва неизползвани дисплеи за изключване. Подробните схеми са в Приложение~В.

\begin{center}
\begin{tikzpicture}[font=\scriptsize,
  stg/.style={draw=navy, fill=navy, text=white, rounded corners=2pt, minimum height=0.85cm, text width=2.55cm, align=center, inner sep=2pt},
  arr/.style={-{Stealth[length=2mm]}, thick, navy}]
\node[stg] (i) {\textbf{Input}\\сензори, код, логове};
\node[stg, right=0.42cm of i] (d) {\textbf{Data Processing}\\MQTT, NAS, псевдонимизация};
\node[stg, right=0.42cm of d] (m) {\textbf{AI Model/Service}\\локални модели на AI сървъра};
\node[stg, right=0.42cm of m] (o) {\textbf{Output}\\прогноза, препоръка};
\node[stg, fill=teal, draw=teal, right=0.42cm of o] (a) {\textbf{Decision/Action}\\одобрение или\\автоматика};
\foreach \x/\y in {i/d,d/m,m/o,o/a} \draw[arr] (\x) -- (\y);
\draw[arr, dashed, teal] (a.south) -- ++(0,-0.3) -| node[pos=0.25, below, text=teal]{Human-in-the-loop: одобренията и корекциите се връщат като нови входни данни} (i.south);
\end{tikzpicture}
\end{center}

\begingroup\footnotesize
\setlength{\tabcolsep}{4pt}
\renewcommand{\arraystretch}{1.1}
\noindent\begin{tabularx}{\textwidth}{@{}>{\bfseries\raggedright\arraybackslash}p{1.95cm}>{\raggedright\arraybackslash}X>{\raggedright\arraybackslash}X>{\raggedright\arraybackslash}X@{}}
\toprule
\rowcolor{sand}Етап & \textbf{AI-1 „СМГ AI тютор“} & \textbf{AI-2 „Умна оранжерия“} & \textbf{AI-3 „График и поддръжка“}\\
\midrule
Input & Въпроси и код на учениците, протоколи; учебни програми, ръководства и задачи на СМГ & Влажност на почвата, температура, светлина, pH, EC; камера; метеостанция & Резервации; часове на печат, цикли на батериите, логове на UPS, консумация; камери на 3D принтерите\\
Data Processing & Вход през Moodle, псевдонимизация, векторен индекс, логове в NAS & MQTT → ThingsBoard; почистване, агрегиране, времеви редове в NAS & Обединяване в Data Hub, нормализиране, маркиране на инциденти\\
AI Model/ Service & Локален отворен LLM (Qwen/Llama чрез Ollama/vLLM) + RAG с цитати & Prophet/LSTM, Isolation Forest, YOLO/MobileNet на Raspberry Pi AI HAT+ / Jetson & Прогноза на търсенето, модел за износване; компютърно зрение за неуспешен печат\\
Output & Подсказки (без готово решение), задачи по ниво, чернова на оценка & Прогноза за 24 часа, сигнал за аномалия, препоръка за поливане & Проект на график, сигнали за поддръжка, дисплеи за изключване\\
Decision/ Action & Учителят одобрява задачите и оценката в Moodle & Автоматично поливане в безопасни граници; торове – с потвърждение от учител & Администраторът одобрява графика; при дефект печатът спира автоматично\\
\bottomrule
\end{tabularx}
\endgroup

\subsection{Рискове от използването на AI и мерки за ограничаването им}

\begingroup\footnotesize
\setlength{\tabcolsep}{4pt}
\renewcommand{\arraystretch}{1.1}
\setlength{\LTpre}{0pt}\setlength{\LTpost}{0pt}
\begin{xltabular}{\textwidth}{@{}>{\bfseries\raggedright\arraybackslash}p{2.55cm}>{\raggedright\arraybackslash}p{5.2cm}>{\raggedright\arraybackslash}X@{}}
\toprule
\rowcolor{sand}Риск & \textbf{Проявление в STEM центъра} & \textbf{Мярка за ограничаване}\\
\midrule
\endfirsthead
\toprule
\rowcolor{sand}Риск & \textbf{Проявление в STEM центъра} & \textbf{Мярка за ограничаване}\\
\midrule
\endhead
Неточност & Грешен отговор на тютора; неточна прогноза за оранжерията & 200 контролни въпроса преди пускане и всеки срок; сравнение на прогнозите с измерванията\\
Халюцинации & Измислени формули, източници или функции от библиотеки & RAG само върху одобрено съдържание на СМГ; задължителен цитат; „Не знам“ при липса на източник\\
Пристрастност (bias) & LLM, обучени главно на английски; малко местни растения в данните & Тестове на български; дообучаване със снимки от оранжерията на СМГ; оценяване по рубрика\\
Защита на данните & Данни на непълнолетни: код, оценки, снимки & Локален модел; псевдонимизация; DPIA; логове до края на годината; камери само към растения и принтери\\
Киберсигурност & Prompt injection; неоторизиран достъп до IoT мрежата и помпата & Роли и права; отделна VLAN за IoT; защита на крайните устройства; ръчен прекъсвач; външен тест (т.~\ref{sec:s3-test})\\
Прекомерна зависимост & Учениците копират отговори, учителите спират да проверяват & Режим „подсказка, не решение“; устна обосновка; без AI на контролни работи; обучение\\
Human-in-the-loop & Оценки, график и дозиране засягат ученици и оборудване & Автоматично – само поливане в граници и спиране на дефектен печат; останалото – с одобрение и лог\\
\bottomrule
\end{xltabular}
\endgroup

\subsection{Външно тестване на AI интеграцията}\label{sec:s3-test}

Приемното тестване на AI интеграцията се възлага на независим външен изпълнител, а не на доставчика по обособената позиция „ИКТ и AI“, за да няма конфликт на интереси; от страна на екипа го координира AI/ML специалистът. То е работен пакет „Интеграция и тестване“ в WBS \lecturetag{Лекция 3: WBS}, а цената и срокът са оценени отдолу нагоре – по тестове и човекодни \lecturetag{Лекция 4: bottom-up оценка}.

\begingroup\footnotesize
\setlength{\tabcolsep}{4pt}
\renewcommand{\arraystretch}{1.1}
\noindent\begin{tabularx}{\textwidth}{@{}>{\raggedright\arraybackslash}X>{\centering\arraybackslash}p{1.8cm}>{\raggedright\arraybackslash}p{4.9cm}@{}}
\toprule
\rowcolor{sand}\textbf{Тест} & \textbf{Човекодни} & \textbf{Критерий за приемане}\\
\midrule
Функционален тест на тютора: 200 контролни въпроса & 3 & ≥\,90\,\% верни отговори\\
Сигурност: prompt injection, изтичане на данни, роли & 3 & 0 критични уязвимости\\
Натоварване: 30 едновременни потребители & 2 & Време за отговор <\,5~s\\
Fail-safe: сензор → AI → поливане при отпадане на сензор & 1 & Поливането спира безопасно\\
Протокол за приемане и повторна проверка & 1 & Подписан протокол\\
\midrule
\textbf{Общо: 10 човекодни × 300~€ + 300~€ инструменти} & \textbf{10} & \textbf{3\,300~€}\\
\bottomrule
\end{tabularx}
\endgroup

\textbf{Време и цена.} Приемното тестване трае три седмици (2–20 август 2027~г.) – след монтажа и близо четири седмици преди въвеждането в експлоатация на 15 септември 2027~г., което оставя резерв за отстраняване на дефекти. От 15 септември до 12 октомври 2027~г. следва четириседмичен пилот с две паралелки; приложение, което не покрие критериите, остава в пилотен режим до повторна проверка. Общо: 3\,300~€ и седем седмици. Сумата се покрива от средствата за организация и управление (\NumMgmt), а пилотът протича в редовните учебни часове без допълнителни разходи.

\clearpage% --- sections/s4.tex
% ============================================================
% Раздел 4. Бизнес и технологичен анализ – SWOT
% ============================================================
\section{Бизнес и технологичен анализ – SWOT}
% етикет в края на абзац: при липса на място се пренася вдясно на нов ред
\newcommand{\sfourtag}[1]{{\unskip\nobreak\hfil\penalty50\hskip1em\hbox{}\nobreak\hfil\mbox{\lecturetag{#1}}\parfillskip=0pt\par}}

SWOT анализът идентифицира и оценява вътрешните (силни и слаби страни) и външните (възможности и заплахи) фактори, които влияят върху изграждането на STEM центъра в СМГ и върху функционирането му след края на финансирането. Анализът е изготвен във фазата на иницииране (ноември–декември 2026~г.) въз основа на оглед на помещенията, инвентаризация на оборудването на Клуб „Роботика“ и лимитите на процедурата; заплахите и слабостите от него са вход за управлението на риска (т.~5). \sfourtag{Лекция 2: SWOT анализ}

\subsection{SWOT матрица}

\begingroup
\footnotesize
\setlength{\tabcolsep}{5pt}
\renewcommand{\arraystretch}{1.0}
\newcommand{\swothead}[2]{\textbf{\color{#1}#2}\par\vspace{1pt}}
\newenvironment{swotlist}{\begin{itemize}[leftmargin=1.0em, itemsep=1pt, topsep=0pt, partopsep=0pt, parsep=0pt, after=\vspace{-0.1\baselineskip}]}{\end{itemize}}
\noindent
\begin{tabularx}{\textwidth}{@{}>{\raggedright\arraybackslash}X!{\color{white}\vrule width 4pt}>{\raggedright\arraybackslash}X@{}}
\rowcolor{navy}
\multicolumn{1}{@{}c!{\color{white}\vrule width 4pt}}{\color{white}\bfseries Положителни фактори} &
\multicolumn{1}{c@{}}{\color{white}\bfseries Отрицателни фактори} \\
\noalign{\vskip 4pt}
\cellcolor{riskL}%
\swothead{teal}{S – Силни страни (вътрешни)}
\begin{swotlist}
  \item \textbf{Интеграция} – 3~направления, 7~учебни предмета и 8~STEM сценария; всеки сензорен поток се ползва от поне 3~предмета.
  \item \textbf{Съвременно оборудване} – ROS2 робот с LiDAR, Jetson Orin Nano, VIS спектрофотометър, GoTo телескоп и AI станция.
  \item \textbf{Локален AI} – 3~AI приложения на сървър в училището; данните на учениците не напускат СМГ.
  \item \textbf{IoT и данни} – „СМГ STEM Data Hub“ (MQTT $\to$ ThingsBoard $\to$ Grafana) свързва 3~лаборатории и 8~ВОСКС.
  \item \textbf{Споделено използване} – лаптопи, дисплеи, 3D принтери и AI сървър по общ график за около 1300~ученици.
  \item \textbf{Човешки капитал} – олимпийска традиция, опитни учители в профилираната подготовка, действащи Клуб „Роботика“ и изследователски клуб.
\end{swotlist}
&
\cellcolor{riskM}%
\swothead{navy}{W – Слаби страни (вътрешни)}
\begin{swotlist}
  \item \textbf{Разходи} – \NumTotal{} (\PctUse{} от лимита) без резерв за поскъпване; след проекта – годишни лицензи (Microsoft 365 Copilot за учителите, MATLAB) и подновяване на ESET след 3~години.
  \item \textbf{Обучение на учителите} – ROS2, IoT и AI са нови за част от колегията; нужно е време извън хорариума, има и съпротива срещу промяната.
  \item \textbf{Техническа сложност} – десетки видове устройства с различни протоколи (USB, Bluetooth, Wi-Fi, MQTT).
  \item \textbf{Интеграция на системите} – ThingsBoard, Moodle/Microsoft~365, AI сървърът и системата за резервации се свързват без външен интегратор; знанието е концентрирано у малко хора.
  \item \textbf{Поддръжка} – калибриране, консумативи, обновления; в училището няма щатен лаборант.
  \item \textbf{Без СМР} – съществуващи помещения (вкл. сутерен) с ограничени площ, вентилация и електрическа мощност.
\end{swotlist}
\\
\noalign{\vskip 4pt}
\cellcolor{sand}%
\swothead{teal}{O – Възможности (външни)}
\begin{swotlist}
  \item \textbf{AI} – по-малки и по-точни отворени модели на същия хардуер; персонализирано обучение и AI в учебните програми.
  \item \textbf{Роботика и интелигентни технологии} – автономни системи, умни сгради, състезания (WRO, FIRST LEGO League).
  \item \textbf{Виртуални лаборатории} – симулации и дигитален близнак на оранжерията: работа извън часовете и с по-малко консумативи.
  \item \textbf{Нови STEM програми} – избираеми и факултативни часове, летни школи, ученически изследователски проекти.
  \item \textbf{Партньорства} – ФМИ на СУ, ТУ-София, БАН и ИТ компании; възпитаници на СМГ в ИТ сектора.
  \item \textbf{Ново финансиране} – Еразъм+, програми на МОН и Столична община; услуги за други училища.
\end{swotlist}
&
\cellcolor{riskH}%
\swothead{navy}{T – Заплахи (външни)}
\begin{swotlist}
  \item \textbf{Технологично остаряване} – лаптопите и GPU остаряват за 4–5 години, AI моделите – за месеци.
  \item \textbf{Киберрискове} – атаки, prompt injection в AI тютора, уязвими IoT устройства, изтичане на лични данни на непълнолетни (GDPR).
  \item \textbf{Недостатъчно използване} – оборудване „в шкафа“ при натиск за подготовка за НВО и ДЗИ.
  \item \textbf{Зависимост от доставчици} – затворени екосистеми (LEGO, NVIDIA, MathWorks, Wolfram), забавени доставки, „най-ниска цена“ по ЗОП.
  \item \textbf{Промени в технологичната среда} – нови стандарти и изискванията на AI Act (Регламент (ЕС) 2024/1689) към AI в образованието.
  \item \textbf{Финансиране след проекта} – несигурни средства от МОН и общината за абонаменти, консумативи и ремонти.
\end{swotlist}
\\
\end{tabularx}
\endgroup

\smallskip
\textbf{Оценка.} Балансът е благоприятен: водещите силни страни са интеграцията и локалният AI, а слабостите са предимно организационни. Най-критични са обучението на учителите (W), киберрисковете и недостатъчното използване (T); те влизат като рискове в RBS и FMEA (т.~5).

\subsection{Стратегически изводи (TOWS)}

Кръстосването на квадрантите превръща SWOT в конкретни действия, които влизат в WBS и в плана на проекта. \sfourtag{Лекция 2: стратегия чрез проекти}

\begin{itemize}[itemsep=2pt]
  \item \textbf{SO – настъпателна (силни страни $\to$ възможности).} Олимпийската традиция, AI сървърът и Data Hub се обединяват в „STEM академия на СМГ“ с ФМИ на СУ, ТУ-София и ИТ компании: поне 8~интердисциплинарни сценария през първата година, менторство от възпитаници, състезания по роботика и отворени дни за други училища.
  \item \textbf{WO – преодоляваща (възможности срещу слабости).} Обучението на място от доставчиците е изискване в техническите спецификации и се допълва с курсове от университетите партньори; ученици от Клуб „Роботика“ стават „STEM асистенти“ – първа линия на поддръжката; виртуалните лаборатории намаляват разхода на реактиви.
  \item \textbf{ST – защитна (силни страни срещу заплахи).} Локален AI модел, псевдонимизация и роли срещу кибер- и GDPR рисковете; отворени стандарти (MQTT, Python, ROS2, отворени LLM) срещу зависимостта от доставчици; AI-3 сигнализира за неизползвано оборудване. AI интеграцията се приема само след \textbf{външно тестване} от независим изпълнител – \textbf{3\,300~€} и \textbf{3~седмици}, последвани от пилот (т.~3.3).
  \item \textbf{WT – отбранителна (минимизиране на слабости и заплахи).} Поръчката е разделена на 3~обособени позиции с гаранционен сервиз в България; безплатен софтуер, където е възможно (GeoGebra, Python/Jupyter, Grafana, Node-RED, Moodle, Ollama); от 2028~г. консумативите и абонаментите се планират в бюджета на училището. Разходите се „ограничават“, обхватът се „подобрява“, а времето се „приема“: при поскъпване се търси нова оферта, а не се превишава лимитът. \sfourtag{Лекция 3: приоритетна матрица}
\end{itemize}

\subsection{Защо интегрираният вариант е по-устойчив}

Алтернативата – три самостоятелни кабинета с отделно оборудване – дублира покупките, ползва всяко устройство само в един предмет и оставя данните разпокъсани. При интегрирания вариант едно оборудване и един поток от данни обслужват 7~предмета и 8~ВОСКС, а поддръжката е обща (една платформа, отворени стандарти, прогнозна поддръжка чрез AI-3). Така всяко вложено евро носи повече учебни часове, а центърът остава полезен и при смяна на учител или доставчик.

\clearpage% --- sections/s5.tex
\section{Управление на риска}

Рисковете са идентифицирани „отдолу нагоре“: за всеки работен пакет от WBS (технически спецификации, обществена поръчка, доставка и монтаж, интеграция на Data Hub, IoT и AI, външно тестване, обучение на учителите, мониторинг) екипът отговаря на въпроса „какво може да се обърка?“ \lecturetag{Лекция 3: WBS}. Всеки риск е отнесен към една от деветте категории на RBS, оценен и възложен на отговорник от проектния екип (раздел 6). Регистърът на рисковете е постоянна точка в месечния отчет на ръководителя на проекта (РП) до директора.

\subsection{Risk Breakdown Structure (RBS)}

RBS има три нива: проект $\to$ 9 категории $\to$ 25 конкретни риска. Всеки риск има код (напр. ФИН-1), който се използва в матрицата и в регистъра. Диаграмата е в Приложение~Г.

\begingroup\footnotesize
\setlength{\tabcolsep}{4pt}
\renewcommand{\arraystretch}{1.12}
\noindent\begin{tabularx}{\textwidth}{@{}>{\raggedright\arraybackslash}p{3.95cm}>{\raggedright\arraybackslash}X@{}}
\toprule
\textbf{Технологични} (ТЕХ) & несъвместимост между доставчиците; повреди; зависимост от един доставчик (vendor lock-in) и остаряване\\
\textbf{Финансови} (ФИН) & поскъпване между концепцията и офертите по ЗОП; недопустими разходи\\
\textbf{Организационни} (ОРГ) & прекратена обособена позиция; забавени доставки; слаба координация\\
\textbf{Образователни} (ОБР) & неготовност на учителите; ниско използване на оборудването в часовете\\
\textbf{AI} & неточност и халюцинации; пристрастност (bias); прекомерна зависимост от AI\\
\textbf{Данни} (ДАН) & загуба на данни в NAS; грешни сензорни данни; изтичане на лични данни\\
\textbf{Киберсигурност} (КИБ) & атака с рансъмуер (ransomware); неоторизиран достъп до MQTT; prompt injection в AI тютора\\
\textbf{Инфраструктурни} (ИНФ) & недостиг на контакти и мощност; слабо покритие на Wi-Fi; влага в сутерена\\
\textbf{Интеграционни} (ИНТ) & сензори към Data Hub; несъвместими платформи; слаба интердисциплинарност\\
\bottomrule
\end{tabularx}
\endgroup

\subsection{Risk Matrix}

Вероятността (P) и въздействието (I) се оценяват по скала 1–5; Risk Score S = P × I. Нива: 1–6 – нисък, 8–12 – среден, 15–25 – висок. От 25-те риска 6 са високи, 17 – средни и 2 – ниски (топлинна карта 5×5 и пълен регистър с ниво, отговорник и мерки за всеки риск – Приложение~Д). Приоритетът следва нивото: мерките за високите рискове се изпълняват още във фазата на планиране (януари–февруари 2027~г.) и се следят на седмичната среща на екипа; средните рискове се преглеждат в месечния отчет; ниските се приемат и се преглеждат всеки срок. Приоритетните рискове (цветът на S показва нивото: червено – висок, жълто – среден):

\begingroup\footnotesize
\setlength{\tabcolsep}{3.5pt}
\renewcommand{\arraystretch}{1.1}
\setlength{\LTpre}{3pt}\setlength{\LTpost}{3pt}
\begin{xltabular}{\textwidth}{@{}>{\bfseries}l>{\raggedright\arraybackslash}p{3.7cm}ccc>{\raggedright\arraybackslash}p{1.8cm}>{\raggedright\arraybackslash}X@{}}
\toprule
\rowcolor{sand}Код & \textbf{Риск} & \textbf{P} & \textbf{I} & \textbf{S} & \textbf{Отговорник} & \textbf{Стратегия: превантивни мерки}\\
\midrule
\endhead
ФИН-1 & Поскъпване на GPU и памети преди ЗОП & 4 & 4 & \cellcolor{riskH}16 & РП, счетоводител & \textit{Намаляване:} пазарни консултации (02.2027); резерв в обхвата (виж по-долу)\\
ОРГ-2 & Забавени доставки (срок 15.09.2027) & 4 & 4 & \cellcolor{riskH}16 & РП & \textit{Прехвърляне:} неустойки за забава; доставка ≤\,8~седмици; резерв 21.08–14.09.2027\\
ОБР-1 & Неготовност и съпротива на учителите & 4 & 4 & \cellcolor{riskH}16 & STEM специалист & \textit{Намаляване:} обучения от доставчиците и вътрешни (06–09.2027); наставник във всяко направление\\
ИНТ-1 & Трудно свързване на сензорите с Data Hub & 4 & 4 & \cellcolor{riskH}16 & IoT инженер & \textit{Намаляване:} MQTT/REST API като изискване в спецификациите; прототип през 02.2027\\
ОРГ-1 & Прекратена обособена позиция по ЗОП & 3 & 5 & \cellcolor{riskH}15 & Юрист по ЗОП & \textit{Намаляване:} спецификации „или еквивалент“; резерв от $\approx$\,3,5~седмици в графика (21.08–14.09.2027) и договаряне без обявление при липса на оферти\\
ДАН-3 & Изтичане на лични данни на непълнолетни & 3 & 5 & \cellcolor{riskH}15 & Експерт по сигурност & \textit{Избягване:} локален AI – данните не напускат училището; DPIA; достъп по роли\\
AI-1 & Халюцинации и грешки на AI тютора & 4 & 3 & \cellcolor{riskM}12 & AI/ML специалист & \textit{Намаляване:} RAG с цитати; 200 контролни въпроса; учителско одобрение; външно тестване\\
КИБ-2 & Неоторизиран достъп до MQTT и помпата & 3 & 4 & \cellcolor{riskM}12 & IoT инженер & \textit{Намаляване:} TLS и удостоверяване на устройствата; ACL по теми; отделна IoT VLAN\\
\bottomrule
\end{xltabular}
\endgroup

Разходите са фиксирани от лимита, затова резервът за ФИН-1 не е парична сума, а предварително подреден списък на позиции, чиито количества се намаляват при поскъпване, без да се засягат SMART целите \lecturetag{Лекция 4: резерв}. Основната мярка за AI-1, КИБ-3 и FMEA режимите F5 и F7 е \textbf{външното тестване на AI интеграцията} от независим изпълнител: 3\,300~€ и 3~седмици през август 2027~г., последвани от пилот (т.~3.3).

\subsection{FMEA на IoT системата на умната оранжерия}

Избран е компонентът, в който AI решение води до физическо действие: сензори $\to$ ESP32 и Raspberry Pi $\to$ MQTT $\to$ ThingsBoard $\to$ AI модел $\to$ помпа и вентилатор. За всеки режим на отказ се оценяват тежест (Impact, I), вероятност (Probability, P) и трудност на откриване (Detection, D) по скала 1–10; Risk Value (RPN) = I × P × D. Критични са отказите с RPN ≥\,150 или I ≥\,9. От 10-те анализирани режима (Приложение~Е) критични са четири:

\begingroup\footnotesize
\setlength{\tabcolsep}{3.5pt}
\renewcommand{\arraystretch}{1.1}
\noindent\begin{tabularx}{\textwidth}{@{}>{\raggedright\arraybackslash}p{4.4cm}cccc>{\raggedright\arraybackslash}X c@{}}
\toprule
\rowcolor{sand}\textbf{Режим на отказ} & \textbf{I} & \textbf{P} & \textbf{D} & \textbf{RPN} & \textbf{Превантивни и коригиращи мерки} & \textbf{RPN след}\\
\midrule
F1 Дрейф на сензора за влажност на почвата & 6 & 6 & 6 & \cellcolor{riskH}216 & Капацитивни сензори; месечно калибриране; контролен сензор и откриване на аномалии & 54\\
F5 Грешна прогноза на AI модела & 6 & 5 & 7 & \cellcolor{riskH}210 & Безопасни граници в контролера извън модела; базово правило за сравнение; преобучение & 72\\
F7 Неоторизиран достъп до MQTT брокера & 9 & 3 & 7 & \cellcolor{riskH}189 & TLS; ACL по теми; IoT VLAN; аларма при нетипични команди & 54\\
F3 Помпата остава включена $\to$ наводнение (I\,=\,9) & 9 & 3 & 5 & \cellcolor{riskH}135 & Хардуерен таймер; датчик за теч; резервоар с ограничен обем; „изключено“ при рестарт & 20\\
\bottomrule
\end{tabularx}
\endgroup

F3 има RPN под 150, но е критичен заради тежестта: наводнение в лабораторията с ток и оборудване е недопустимо дори при малка вероятност. Общото правило е fail-safe: при загуба на връзка или захранване и при съмнителна прогноза контролерът преминава към локални правила с безопасни граници, а дозирането на торове изисква потвърждение от учител (human-in-the-loop). След мерките сумарният RPN на десетте режима спада от 1\,273 на 350 (−72,5\,\%). Веригата сензор $\to$ AI $\to$ поливане е част от външното тестване през август 2027~г.; автоматичното поливане се включва само след успешен тест за безопасно спиране.

\clearpage% --- sections/s6.tex
\section{Управление на проектния екип}

Проектът се изпълнява от екип на СМГ, организиран като \textbf{слаба матрична структура}: членовете остават на основните си длъжности и отделят част от работното си време (ръководителят на проекта – около 30\,\%, ключовите специалисти – 10–20\,\%). Ръководител на проекта (РП) е заместник-директорът по учебната дейност, тъй като той разпределя часовете, кабинетите и натовареността на учителите – ресурсите, без които центърът не може да функционира. Ролята му е интегрираща, а не техническа: обединява решенията на специалистите в общ план, поддържа баланса в триъгълника обхват–време–разходи и мотивира екипа \mbox{\lecturetag{Лекция 1: роля на ръководителя}}. Под РП работят три блока – образователен, „AI и софтуер“ и „Инфраструктура, IoT и роботика“. Екипът обединява 15 служители на училището, двама външни експерти и ученици от Клуб „Роботика“; организационната схема и матрицата RACI по работните пакети от WBS са в Приложение~Ж.

\subsection{Роли, отговорности и компетентности}

{\footnotesize
\setlength{\tabcolsep}{4pt}
\renewcommand{\arraystretch}{1.08}
\begin{tabularx}{\textwidth}{@{}>{\raggedright\arraybackslash}p{4.3cm}>{\raggedright\arraybackslash}X>{\raggedright\arraybackslash}p{4.4cm}@{}}
\toprule
\textbf{Роля – изпълнител} & \textbf{Основни отговорности} & \textbf{Компетентности}\\
\midrule
\textbf{Спонсор} – директор & Одобрява устава на проекта, бюджета и договорите; решава ескалираните спорове & Стратегическа визия; правомощия на възложител\\
\textbf{Project Manager} – зам.-директор УД & План, WBS, график, ресурси, контрол на сроковете, отчети & Управление на проекти, лидерство, комуникация\\
\textbf{STEM Integration Specialist} – учител по физика & STEM Integration Matrix, сценарии, общ график на лабораториите & Интердисциплинарна методика, проектно обучение\\
\textbf{AI/ML Specialist} – учител по информатика + външен консултант & AI тютор (AI-1), модели за оранжерията (AI-2) и за графика (AI-3); координира външното тестване на AI & LLM и RAG, машинно обучение, оценка на модели, етика на AI\\
\textbf{Software Developer} – учител по ИТ + ученици от Клуб „Роботика“ (под надзор) & „СМГ STEM Data Hub“: ThingsBoard, Grafana, резервации, връзка с Moodle/Microsoft~365 & Python, бази данни, API, Git\\
\textbf{IoT/Robotics Engineer} – ръководител на Клуб „Роботика“ & Сензори, MQTT, оранжерия, роботи и дронове, безопасни граници на автоматиката & Електроника, микроконтролери, ROS2, безопасност\\
\textbf{Infrastructure Specialist} – зам.-директор АСД + системен администратор & Помещения, електрически товар, UPS, мрежа, сървър и NAS, инвентар & Мрежи, сървъри, управление на материалната база\\
\textbf{Cybersecurity Specialist} – външен експерт по договор & Модел на заплахите, достъп по роли, защита на данните (GDPR), критерии за теста за сигурност & Мрежова и AI сигурност (prompt injection)\\
\textbf{Teacher/Subject Expert} – учители по математика, биология, химия, география & Изисквания по предмет, спецификации, уроци по сценариите, пилот & Предметна експертиза, дигитална педагогика\\
\textbf{Data Analyst} – учител по математика & KPI, данни за използването, STEM Integration Score & Статистика, визуализация (Grafana, Python)\\
\textbf{ЗОП и финанси} – юрист по ЗОП, счетоводител & Документация по ЗОП (3 обособени позиции), договори, плащания, отчитане & ЗОП, правила на НПВУ, финансов контрол\\
\bottomrule
\end{tabularx}}

\subsection{Развитие на екипа по модела на Тъкман}

Етапите следват жизнения цикъл; стилът на РП преминава от насочване към делегиране.

{\footnotesize
\setlength{\tabcolsep}{4pt}
\renewcommand{\arraystretch}{1.08}
\begin{xltabular}{\textwidth}{@{}>{\raggedright\arraybackslash}p{1.85cm}>{\raggedright\arraybackslash}p{5.2cm}>{\raggedright\arraybackslash}X@{}}
\toprule
\textbf{Етап} & \textbf{Очаквани прояви и конфликти} & \textbf{Управленски действия}\\
\midrule
\endhead
\textbf{Forming}\newline ноем.–дек. 2026 & Ентусиазъм, но неясни роли; скептицизъм на част от учителите към AI и към „поредния проект“ & Встъпителна среща, устав на проекта, SMART цели, RACI; декларации за участие на учителите и решение на Педагогическия съвет; стил „насочване“\\
\textbf{Storming}\newline ян.–февр. 2027 & Марки и платформи (LEGO SPIKE Prime или Arduino/ESP32); дял на бюджета по направления; часове в графика на лабораториите; политика за AI („забрана“ или „разрешаване“) & Решения по критерии (брой ползващи предмети, съвместимост с Data Hub, цена); функционални спецификации („или еквивалент“ по ЗОП); при фиксирани разходи обхватът се договаря по приоритетната матрица; политика за AI, обсъдена с учители, ученици и родители; ескалация до директора след 7 дни; стил „обучаване“\\
\textbf{Norming}\newline март–май 2027 & Приети правила и доверие; риск от „групово мислене“ при оценката на офертите & Работни правила (отговор до 48 часа, общи шаблони); „адвокат на дявола“ и декларации за безпристрастност при оценката на офертите; стил „подкрепа“\\
\textbf{Performing}\newline юни 2027 – юни 2028 & Самостоятелни подекипи по обособените позиции; натиск от забавени доставки и от интеграция през ваканцията & Делегиране и контрол на критичния път; график на отпуските за август; външно тестване на AI (2–20 август 2027~г., 3\,300~€) и 4-седмичен пилот; при сериозно забавяне – ретроспектива (временно връщане към Storming)\\
\textbf{Adjourning}\newline юли 2028 & Тревога за поддръжката след проекта; спад на мотивацията & Финален отчет, STEM Integration Score, извлечени поуки, признание; поддръжката се поема от постоянен STEM съвет\\
\bottomrule
\end{xltabular}}

\subsection{Комуникация и координация}

Координацията следва комуникационния план \lecturetag{Лекция 3: комуникационен план}: ежеседмична 30-минутна оперативна среща на екипа в сряда от 14:30~ч. – „какво свърших, какво предстои, какво ме блокира“; месечен писмен отчет на РП до директора (напредък по WBS, бюджет, рискове, светофар зелено/жълто/червено); общ канал в Microsoft Teams с папка за всеки работен пакет – единен архив на документите и решенията; отчет пред Педагогическия съвет всеки учебен срок. Външното приемно тестване на AI интеграцията (т.~3.3) се координира от AI/ML специалиста (R), РП се отчита за него пред спонсора (A), а специалистът по киберсигурност и анализаторът на данни задават критериите за приемане.

\clearpage% --- sections/s7.tex
\section{Използване на изкуствен интелект при управлението на проекта}

Подготовката на проекта е проведена по модела \emph{AI-Augmented PM} \lecturetag{Лекция 1: Traditional → Digital → AI-Augmented PM}. Генеративният AI работи като асистент, а решенията остават за ръководителя на проекта. Използван е Claude Code (Anthropic, модел Claude Opus~5.5) в работни потоци с паралелни агенти: един агент изпълнява задачата, а друг, независим агент проверява резултата. Решенията, които AI не може да вземе сам (кои направления, СМР или не, как се третира ДДС, кои лимити важат), бяха зададени като въпроси към студента и записани в регистъра.

\subsection{Регистър на използването на AI (AI Project Log)}
Пълният регистър с 12 записа е в Приложение~З (колони: задача, AI инструмент, заявка, генериран резултат, приет/отхвърлен, проверка). Основните записи са:
\begin{itemize}
  \item \textbf{Проучване на цените} – 10 изследователски и 10 проверяващи агента обработиха 106 позиции, плюс 2 агента за 3 допълнителни позиции за ВОСКС. \emph{Резултат:} 109 проучени позиции (81 при двама български търговци и 28 по оферта); в бюджета остават \NumItems{} (\NumBGItems{} + \NumOffers): 5 предложения са отхвърлени (виж по-долу) и един модел OPS компютър е обединен. \emph{Проверка:} независимият агент зареди отново всяка страница. 88 позиции бяха потвърдени без промяна, а 21 (19\,\%) бяха коригирани: изчерпани продукти, грешна марка по EAN, неточни спецификации, оферта при наличие на два български търговеца, два магазина от една търговска мрежа.
  \item \textbf{Чернови на раздели 1–6 и 8} – 7 пишещи и 7 рецензиращи агента. \emph{Проверка:} рецензентите откриха несъответствия между разделите (различен ръководител на проекта в т.~1 и т.~6, 74 вместо 76 места, логери в грешната лаборатория) и ги поправиха; всички изчисления в таблиците бяха преизчислени.
  \item \textbf{Отхвърлени предложения} – виртуални лаборатории (Labster, около 8\,600~€/год.) и хостван инвентар (Snipe-IT), защото софтуерът щеше да надхвърли 10\,\%; компактният AI компютър (ASUS GX10, около 10\,500~€) и професионалната термокамера (FLIR E8 Pro), защото специализираното оборудване щеше да надхвърли 27\,000~€ (за GX10 няма и текущ български търговец); офертата за NetSupport от САЩ – заменена с български търговец.
  \item \textbf{Скриптове за бюджета} – AI генерира програма, която попълва образеца и оптимизира количествата. Ръчната проверка на сумите (J45 $\neq$ J43 + J44) откри грешка в генерирания изчислител на формули; тя беше поправена, преди да бъдат използвани числата.
\end{itemize}

\subsection{Сравнение: Human-only срещу AI-assisted}
За AI-assisted подхода времето е измерено от логовете на работните потоци. За Human-only подхода то е оценено „отдолу нагоре“ \lecturetag{Лекция 4: оценка}: около 5~минути за намиране и проверка на една цена, 15~минути за една оферта и около 3~часа за раздел.

{\footnotesize
\begin{xltabular}{\textwidth}{@{}>{\raggedright\arraybackslash}p{2.5cm} >{\raggedright\arraybackslash}X >{\raggedright\arraybackslash}X >{\raggedright\arraybackslash}X >{\raggedright\arraybackslash}X@{}}
\toprule
\textbf{Показател} & \textbf{Цени – Human-only} & \textbf{Цени – AI-assisted} & \textbf{Раздели – Human-only} & \textbf{Раздели – AI-assisted} \\
\midrule
Време & $\approx$ 22~ч (81$\times$2$\times$5~мин + 28$\times$15~мин + 2~ч Excel) & 48 + 24 + 19~мин агенти; + около 2~ч за решения и преглед & $\approx$ 21~ч (7 раздела $\times$ 3~ч) & 18,5~мин; + около 3~ч за преглед и редакция \\
Качество & 1 вариант на позиция, без проверка от втори човек & 109/109 проверени от втори агент; методиката и източниците са в листа & зависи от автора, рядко се проверява съгласуваност & 7 независими рецензии; общ лист с факти за всички раздели \\
Точност & рискът от грешна цена или изчерпан продукт не е измерен & 21 корекции (19\,\%) преди окончателния вариант; 0 неуспешни проверки & типични грешки при пресмятанията & несъответствията между разделите са открити и поправени; аритметиката е проверена \\
Алтернативи & 1 & за част от позициите са сравнени 2–4 модела (напр. 4 системи за безжично представяне) & 2–3 варианта & 4 варианта в т.~8, 4 TOWS стратегии в т.~4 \\
\bottomrule
\end{xltabular}}

\subsection{Изводи}
AI съкрати времето за двете задачи около 6 пъти с включения преглед (само работата на агентите – над 10 пъти) и даде повече алтернативи. Без проверка обаче около всяка пета цена щеше да е неточна: AI генерира бързо, но проверката трябва да е заложена в процеса като отделна стъпка с независим проверяващ. Решенията с последици за обхвата и бюджета останаха за ръководителя на проекта (human-in-the-loop). Разходът на ресурси също е измерим: около 6,3~млн. токена за всички агенти (4,5~млн. за цените, вкл. ВОСКС и повторната проверка, и 1,8~млн. за разделите). При тарифиране по API цената е (брой токени $\div$ 1~млн.) $\times$ цената за 1~млн. токена, т.е. $\approx$\,6,3 $\times$ тарифата; при абонамент допълнителен разход няма. Това е същата логика като при външното тестване на AI интеграцията в раздел~3: AI се приема само след проверка, за която са предвидени време и разходи.

\clearpage% --- sections/s8.tex
\section{Оптимизация на STEM интеграцията}

Целта е да се избере вариантът с най-висока степен на интеграция при един и същ фиксиран лимит. Според приоритетната матрица на проекта разходите са „ограничени“ (300\,000~€ за STEM и 60\,000~€ за ВОСКС), времето се „приема“, а обхватът се „подобрява“~– следователно въпросът не е кой вариант е най-евтин, а кой постига най-пълна интеграция в рамките на бюджета \lecturetag{Лекция 3: приоритетна матрица}.

\subsection{Модел и разглеждани варианти}

Използва се претеглен многокритериален модел с девет критерия, всеки от които се оценява по петстепенна скала (1~– много ниска, 2~– ниска, 3~– средна, 4~– висока, 5~– много висока степен на изпълнение):
\[
  \textstyle \text{STEM Integration Score} = \sum_{i=1}^{9} W_i \times X_i, \qquad \sum_{i=1}^{9} W_i = 1, \qquad
  \text{Score}_{\%} = \dfrac{\text{Score}}{5} \times 100\,\%.
\]
Критериите „Разходи за реализация“ и „Сложност на управлението“ са с обърната скала (5 = най-изгодно, съответно най-просто), така че по-високата оценка винаги е по-добра. Четирите критерия, които изразяват същността на интеграцията (дисциплини, споделени технологии, информационни потоци, AI), носят общо 0,65 от теглото, защото именно те са целта на процедурата; използването на оборудването (0,10) и разширяемостта (0,08) определят дългосрочната полза, а разходите получават само 0,07, тъй като лимитът е фиксиран и всички варианти се вписват в него. Процентът от максимума показва дял от възможните точки по модела, а не непременно „реален“ процент интеграция.

Сравняват се четири варианта: \textbf{A}~– „Три отделни кабинета“ (класически, без общ хъб и без AI); \textbf{B}~– „Облачно-центриран“ (SaaS AI и облачна IoT платформа, без локален сървър); \textbf{C}~– „Интегриран център със СМГ STEM Data Hub и локален AI“ (три лаборатории и 8 ВОСКС, обединени от локален сървър, ThingsBoard/MQTT, NAS и три AI приложения); \textbf{D}~– „Един голям многофункционален кабинет“ (всички направления в обща зала с общ комплект оборудване).

\subsection{Многокритериална матрица за оценка и анализ на чувствителността}

{\footnotesize
\setlength{\tabcolsep}{3.5pt}
\begin{tabularx}{\textwidth}{@{}X c cc cc >{\columncolor{sand}}c >{\columncolor{sand}}c cc@{}}
\toprule
 & & \multicolumn{2}{c}{\textbf{A}} & \multicolumn{2}{c}{\textbf{B}} & \multicolumn{2}{c}{\cellcolor{sand}\textbf{C (избран)}} & \multicolumn{2}{c@{}}{\textbf{D}} \\
\cmidrule(lr){3-4}\cmidrule(lr){5-6}\cmidrule(lr){7-8}\cmidrule(l){9-10}
\textbf{Критерий} & $\mathbf{W_i}$ & $X_i$ & $W_iX_i$ & $X_i$ & $W_iX_i$ & $X_i$ & $W_iX_i$ & $X_i$ & $W_iX_i$ \\
\midrule
1. Интеграция на учебни дисциплини & 0,20 & 2 & 0,40 & 4 & 0,80 & 5 & 1,00 & 3 & 0,60 \\
2. Споделено използване на технологии и оборудване & 0,15 & 2 & 0,30 & 3 & 0,45 & 5 & 0,75 & 4 & 0,60 \\
3. Взаимосвързани информационни потоци & 0,15 & 1 & 0,15 & 4 & 0,60 & 5 & 0,75 & 2 & 0,30 \\
4. AI интеграция & 0,15 & 1 & 0,15 & 4 & 0,60 & 5 & 0,75 & 2 & 0,30 \\
5. Степен на използване на оборудването & 0,10 & 2 & 0,20 & 3 & 0,30 & 4 & 0,40 & 4 & 0,40 \\
6. Възможности за разширяване & 0,08 & 3 & 0,24 & 3 & 0,24 & 5 & 0,40 & 2 & 0,16 \\
7. Разходи за реализация (5 = най-изгодно) & 0,07 & 4 & 0,28 & 3 & 0,21 & 3 & 0,21 & 4 & 0,28 \\
8. Енергийна ефективност & 0,05 & 3 & 0,15 & 4 & 0,20 & 3 & 0,15 & 4 & 0,20 \\
9. Сложност на управлението (5 = най-просто) & 0,05 & 4 & 0,20 & 4 & 0,20 & 3 & 0,15 & 4 & 0,20 \\
\midrule
\textbf{STEM Integration Score} & \textbf{1,00} & & \textbf{2,07} & & \textbf{3,60} & & \textbf{4,56} & & \textbf{3,04} \\
\textbf{Процент от максимума (Score/5)} & & & 41,4\,\% & & 72,0\,\% & & \textbf{91,2\,\%} & & 60,8\,\% \\
\midrule
Чувствителност: равни тегла $W_i = 1/9$ & & & 2,44 & & 3,56 & & \textbf{4,22} & & 3,22 \\
Чувствителност: разходи 0,20, дисциплини 0,07 & & & 2,33 & & 3,47 & & \textbf{4,30} & & 3,17 \\
\bottomrule
\end{tabularx}
}

Класирането C > B > D > A се запазва и при двата алтернативни сценария на теглата: C получава 84,4\,\%, съответно 86,0\,\% от максимума, а преднината му пред B е 0,66 и 0,83. Дори при пълно прехвърляне на теглото на критерий~1 към разходите (0 и 0,27) C остава първи с 4,16 срещу 3,40 за B, т.\,е. изборът е устойчив спрямо субективната преценка при определянето на теглата.

\subsection{Обосновка на оценките}

\textbf{A} запазва модела „всеки предмет със свое оборудване“: данните не се обменят (потоци~– 1), няма AI (1), а оборудването стои неизползвано извън часовете по съответния предмет (2). Вариантът е евтин и лесен за управление (по 4), но не отговаря на целта на процедурата. \textbf{D} използва оборудването интензивно (4) и е прост за управление (4), но в една зала може да работи само една паралелка наведнъж, няма общ хъб за данни, AI е ограничен (потоци и AI~– 2), а мястото за нови направления е изчерпано (разширяване~– 2).

\textbf{B} е добре свързан (потоци и AI~– 4), но данните на непълнолетните ученици се обработват на външни сървъри, често извън ЕС, което изисква договор по GDPR; доставчикът обновява модела без предупреждение и приемното тестване не е възпроизводимо. Инвестицията е по-ниска, но абонаментите остават и след края на финансирането (разходи~– 3).

\textbf{C} получава 5 по четирите интеграционни критерия: 7 предмета и 8 интердисциплинарни сценария; лаптопите, дисплеите, 3D принтерите и AI сървърът се ползват от трите направления и 8-те ВОСКС по общ график; веригата сензор → MQTT → ThingsBoard → NAS → AI → табло → решение свързва всички помещения, а данните на учениците не напускат училището. Оценката 5 за AI предполага и \textbf{външно приемно тестване} на интеграцията (т.~3.3). Разходите са 3: най-високата инвестиция (AI станция, NAS, специализирано оборудване) остава в рамките на лимита, а отворените модели не изискват лицензи. Енергийната ефективност и сложността (по 3) отразяват постоянно работещия сървър и IoT мрежата и се смекчават от AI-3 и екипа по т.~6.

\subsection{Извод}

Оптималният вариант е \textbf{C} със STEM Integration Score 4,56 от 5 (91,2\,\% от максимума по модела) и преднина 0,96 пред B. Той използва \PctUse{} от общия лимит и превръща фиксирания бюджет в максимален обхват на интеграцията. При закриването (юли 2028~г.) моделът ще се приложи повторно с реални данни за използването \lecturetag{Лекция 1: жизнен цикъл}.

\clearpage% --- sections/s9.tex
\section{Финансов анализ}

Изходният резултат по т.~9 е Excel файлът „Концепция за изграждане на училищна STEM среда“~– попълненият официален образец с листове за трите направления („Направление 5, 1 и 3“), „ВОСКС“, „Специализирано оборудване“ и добавен лист „Методика и източници“.

\subsection{Методика: оценка „отдолу нагоре“ с проверка „отгоре надолу“}
Бюджетът е изграден по метода \emph{bottom-up} \lecturetag{Лекция 4: bottom-up и top-down}: за всяка от \NumItems{} позиции (\NumLines{} реда в листовете) е избран конкретен модел и е определена единична цена. За \NumBGItems{} позиции моделът е намерен при \textbf{двама различни български търговци} (общо \NumSellers{} търговци); в листа са линковете и двете цени с ДДС, колона~N е средната цена, а единичната цена без ДДС е $N \div 1{,}2$. За \NumOffers{} позиции, които не се продават от български търговец (\NumOffersEU{} от ЕС и \NumOffersNonEU{} извън ЕС), е използвана \textbf{„Оферта \#N“}: (цена в чужбина + допълнителни разходи) $\times$ 1,35. Допълнителните разходи следват правилата за внос: от ЕС~– без мито и без допълнително ДДС; извън ЕС~– 20\,\% ДДС над 15,34~€, мито по тарифа (поне 4\,\%) и такса за обработка 20~€ над 150~€ (3~€ за малки пратки); софтуерът от трети страни се облага само с ДДС. Изчислението е в колони O–AB на всеки ред. Цените са проверени на 05.11.2026~г. (курсове на ЕЦБ от 04.11.2026~г.), като всеки източник е зареден повторно от независим проверяващ агент (раздел~7).

Проверката \emph{top-down} е сравнението с лимитите (табл.~9.1): количествата на гъвкавите позиции (напр. сензорни логери, комплекти LEGO, дронове, VR очила, нишки за 3D печат) са изчислени така, че сумите да са възможно най-близо до лимитите, без да ги надвишават.

% --- sections/s9_tables.tex
\begin{table}[h]\centering\footnotesize
\begin{tabularx}{\textwidth}{@{}l X r r r l@{}}
\toprule
\textbf{№} & \textbf{Вид разход} & \textbf{Лимит / праг} & \textbf{Искано} & \textbf{\% от лимита} & \textbf{Правило} \\
\midrule
3.1 & Средства STEM & 300\,000,00~€ & 299\,998,58~€ & 99,9\,\% & ≤ лимита \textcolor{teal}{$\checkmark$} \\
3.2 & СМР и обзавеждане (СМР – 0~€) & 135\,000,00~€ & 36\,039,17~€ & 26,6\,\% & ≤ 45\,\% \textcolor{teal}{$\checkmark$} \\
3.3 & Оборудване (3 направления) & мин. 135\,000,00~€ & 211\,934,26~€ & 156,9\,\% & ≥ 45\,\% \textcolor{teal}{$\checkmark$} \\
3.4 & Допълнително специализирано & 27\,000,00~€ & 22\,258,45~€ & 82,4\,\% & ≤ 20\,\% от 3.3 \textcolor{teal}{$\checkmark$} \\
3.5 & Софтуер STEM & 30\,000,00~€ & 29\,766,71~€ & 99,2\,\% & ≤ 10\,\% \textcolor{teal}{$\checkmark$} \\
3.6 & Средства ВОСКС & 60\,000,00~€ & 59\,984,12~€ & 99,9\,\% & ≤ лимита \textcolor{teal}{$\checkmark$} \\
3.7 & в т.ч. софтуер ВОСКС & 6\,000,00~€ & 1\,219,20~€ & 20,3\,\% & ≤ 10\,\% \textcolor{teal}{$\checkmark$} \\
3.8 & Общо STEM + ВОСКС & 360\,000,00~€ & 359\,982,71~€ & 99,9\,\% & ≤ лимита \textcolor{teal}{$\checkmark$} \\
3.9 & Организация и управление (1\,\%) & 3\,600,00~€ & 3\,599,83~€ & 99,9\,\% & = 1\,\% от 3.8 \textcolor{teal}{$\checkmark$} \\
3.10 & \textbf{Общо} & 363\,600,00~€ & 363\,582,54~€ & 99,9\,\% & ≤ лимита \textcolor{teal}{$\checkmark$} \\
\bottomrule
\end{tabularx}
\caption*{\footnotesize Таблица 9.1. Искано финансиране спрямо лимитите (€ с ДДС). Неизползван остатък: \NumLeftTotal.}
\end{table}
\begin{table}[h]\centering\footnotesize
\begin{tabularx}{\textwidth}{@{}X r r r r r@{}}
\toprule
\textbf{Лист} & \textbf{Оборудване} & \textbf{Софтуер} & \textbf{Обзавеждане} & \textbf{Общо} & \textbf{Дял} \\
\midrule
Математика и информатика & 63\,632,32 & 20\,625,56 & 4\,350,82 & 88\,608,70 & 24,6\,\% \\
Роботика и КФС & 73\,806,22 & 8\,625,24 & 9\,097,04 & 91\,528,50 & 25,4\,\% \\
Природни науки & 74\,495,72 & 515,90 & 22\,591,31 & 97\,602,94 & 27,1\,\% \\
Специализирано оборудване & 22\,258,45 & 0,00 & – & 22\,258,45 & 6,1\,\% \\
ВОСКС (8 класни стаи) & 58\,764,92 & 1\,219,20 & – & 59\,984,12 & 16,6\,\% \\
\midrule
\textbf{Общо STEM + ВОСКС} & 292\,957,63 & 30\,985,91 & 36\,039,17 & \textbf{359\,982,71} & 100\,\% \\
\bottomrule
\end{tabularx}
\caption*{\footnotesize Таблица 9.2. Разпределение по листове на образеца (€ с ДДС).}
\end{table}


\subsection{Структура и изводи}
Използвани са \PctUse{} от общия лимит~– неизползваният остатък е само \NumLeftTotal. Оборудването (\NumEquip) е далеч над минимума от 45\,\%~– приоритет е работата на учениците, а не обзавеждането. Обзавеждането (\NumFurn) е далеч под тавана от 45\,\%, защото не се правят СМР, а помещенията се обновяват с лабораторни острови, работни маси и акустични панели. Софтуерът (\NumSoft) почти изчерпва 10\,\%: ThingsBoard и Mathematica са с безсрочни лицензи, ESET и MATLAB~– за 3~години, а годишни са само лицензите за Microsoft 365 Copilot~– за учителите, без лични данни на ученици. Най-големите редове са: \TopLines. По оферти са \PctOfferTotal{} от бюджета (\NumOfferTotal), главно лабораторни уреди от доставчици в ЕС.

\subsection{Резерв, устойчивост и разходи извън проекта}
\textbf{Резерв.} Образецът не позволява ред „непредвидени разходи“, а лимитите са фиксирани. Затова резервът е „вграден“ в количествата: \NumFlexBuffer{} са в бройки над базовото ниво на гъвкавите позиции. Ако цените при поръчката по ЗОП се окажат по-високи, тези количества се намаляват, без да се слиза под минималното ниво по направления. Така разходите остават „ограничени“, а обхватът се „подобрява“ или се свива контролирано \lecturetag{Лекция 3: приоритетна матрица}.

\textbf{Външно тестване на AI интеграцията.} От средствата за организация и управление се плащат 3\,300~€ за независимото приемно тестване (10~човекодни, 3~седмици през август 2027~г.; т.~3.3); останалото покрива административни разходи, а трудът на училищния екип не се финансира от безвъзмездните средства.

\textbf{Годишни разходи след проекта} (бюджет на училището, ориентировъчно): консумативи за 3D печат, химия и електроника~– около 2\,500~€; годишните лицензи (Microsoft 365 Copilot за учителите и MATLAB)~– около 3\,000~€; подновяване на ESET след 3~години~– средно около 3\,000~€ годишно; резервни части и поддръжка~– около 1\,500~€. Общо около 10\,000~€ годишно, т.е. под 3\,\% от инвестицията. Безсрочните лицензи и отвореният софтуер ограничават зависимостта от абонаменти.

% --- sections/appendix.tex
\clearpage
\section*{Приложения}
\addcontentsline{toc}{section}{Приложения}
Приложенията съдържат пълните таблици и диаграми, към които препращат раздели 2, 3, 5, 6 и 7. Бюджетът (раздел 9) е отделен файл във формат Excel – попълненият образец на Концепция за изграждане на училищна STEM среда.

% --- sections/a_s2.tex
% ---------------- Приложения към т. 2 ----------------
% Цветове и значки на учебните предмети (общи за Приложение А и Б)
\providecolor{sIIcM}{HTML}{1F3A5F}  % Математика
\providecolor{sIIcI}{HTML}{2A7F7F}  % Информатика
\providecolor{sIIcIT}{HTML}{6A4C93} % Информационни технологии
\providecolor{sIIcF}{HTML}{B03A2E}  % Физика и астрономия
\providecolor{sIIcH}{HTML}{C77C0E}  % Химия и ООС
\providecolor{sIIcB}{HTML}{2E7D4F}  % Биология и ЗО
\providecolor{sIIcG}{HTML}{8B5A2B}  % География и икономика
\providecommand{\sIIbadge}[2]{{\setlength{\fboxsep}{0.8pt}\colorbox{#1}{\makebox[8.6pt]{\textcolor{white}{\bfseries\tiny #2}}}}}
\providecommand{\sIIbM}{\sIIbadge{sIIcM}{М}}
\providecommand{\sIIbI}{\sIIbadge{sIIcI}{И}}
\providecommand{\sIIbIT}{\sIIbadge{sIIcIT}{ИТ}}
\providecommand{\sIIbF}{\sIIbadge{sIIcF}{Ф}}
\providecommand{\sIIbH}{\sIIbadge{sIIcH}{Х}}
\providecommand{\sIIbB}{\sIIbadge{sIIcB}{Б}}
\providecommand{\sIIbG}{\sIIbadge{sIIcG}{Г}}
\providecommand{\sIIball}{\sIIbM\hspace{0.5pt}\sIIbI\hspace{0.5pt}\sIIbIT\hspace{0.5pt}\sIIbF\hspace{0.5pt}\sIIbH\hspace{0.5pt}\sIIbB\hspace{0.5pt}\sIIbG}

\subsection*{Приложение А. STEM Integration Matrix}
\addcontentsline{toc}{subsection}{Приложение А. STEM Integration Matrix}

Всеки ред е верига \emph{учебна дисциплина → STEM област → технология → оборудване → данни → AI → учебна дейност} (към т.~2). Под водещата дисциплина са отбелязани и другите предмети, които използват същата верига: \sIIbM~математика, \sIIbI~информатика, \sIIbIT~информационни технологии, \sIIbF~физика и астрономия, \sIIbH~химия и ООС, \sIIbB~биология и ЗО, \sIIbG~география и икономика. Пространства: Л1 – лаборатория „Данни и AI“; Л2 – „Роботика и мейкърспейс“; Л3 – „Природни науки“ с умната оранжерия; DH – „СМГ STEM Data Hub“; ВОСКС – 8 високотехнологични класни стаи; (спец.) – допълнително специализирано оборудване.

{\scriptsize
\setlength{\tabcolsep}{3pt}
\renewcommand{\arraystretch}{1.12}
\begin{xltabular}{\textwidth}{@{}>{\raggedright\arraybackslash}p{0.34cm}
  >{\hsize=0.95\hsize\raggedright\arraybackslash\hspace{0pt}}X
  >{\hsize=0.8\hsize\raggedright\arraybackslash\hspace{0pt}}X
  >{\hsize=0.85\hsize\raggedright\arraybackslash\hspace{0pt}}X
  >{\hsize=1.2\hsize\raggedright\arraybackslash\hspace{0pt}}X
  >{\hsize=1.0\hsize\raggedright\arraybackslash\hspace{0pt}}X
  >{\hsize=1.0\hsize\raggedright\arraybackslash\hspace{0pt}}X
  >{\hsize=1.2\hsize\raggedright\arraybackslash\hspace{0pt}}X@{}}
\toprule
\textbf{№} & \textbf{Дисциплина} & \textbf{STEM област} & \textbf{Технология} & \textbf{Оборудване (пространство)} & \textbf{Данни} & \textbf{AI} & \textbf{Учебна дейност} \\
\midrule
\endfirsthead
\toprule
\textbf{№} & \textbf{Дисциплина} & \textbf{STEM област} & \textbf{Технология} & \textbf{Оборудване (пространство)} & \textbf{Данни} & \textbf{AI} & \textbf{Учебна дейност} \\
\midrule
\endhead
\bottomrule
\endfoot
\multicolumn{8}{@{}l}{\cellcolor{sand}\textbf{Направление „Природни науки“ – Л3 и умна оранжерия}}\\
1 & \textbf{Биология}\newline \sIIbH\,\sIIbF\,\sIIbM\,\sIIbI & Физиология и екология на растенията & IoT сензори, компютърно зрение & Умна оранжерия: сензори за почва и светлина; камера (Л3) & Температура, влажност, светлина (на 1~min); снимки & AI-2: прогноза на времеви редове (Prophet/LSTM), YOLO/MobileNet & „Умна оранжерия“: прогноза на условията за растеж, план за поливане \\
2 & \textbf{Химия и ООС}\newline \sIIbB\,\sIIbM & Аналитична химия & Електрохимични сензори, VIS спектрофотометрия (335–800~nm) & pH метри, EC сензор, логери, спектрофотометър (спец.) (Л3) & pH, EC, спектри на хранителния разтвор & AI-2: аномалии (Isolation Forest); AI-1: обратна връзка по протокола & Контрол на разтвора; дозиране на торове само след потвърждение от учител \\
3 & \textbf{Биология}\newline \sIIbI\,\sIIbH & Клетъчна биология & Цифрова микроскопия, компютърно зрение & 14 цифрови микроскопа, тринокулярен микроскоп с камера (Л3); AI сървър (DH) & Микроснимки на клетки и тъкани, анотирани от учениците & AI-2 (лабораторни данни): CV класификатор & „AI под микроскоп“: класификация на клетки, сравнение човек – модел \\
4 & \textbf{Физика и астрономия}\newline \sIIbM\,\sIIbI & Астрономия и оптика & Астро\-фото\-гра\-фия, обработка на изображения & GoTo телескоп с астрокамера (спец.) (Л3); AI сървър (DH) & Изображения FITS, времеви серии на яркостта & Учебен CV модел за откриване на обекти; статистика на шума & „Нощното небе“: криви на блясъка, размери на кратери \\
\multicolumn{8}{@{}l}{\cellcolor{sand}\textbf{Направление „Роботика и кибер-физични системи“ – Л2}}\\
5 & \textbf{Физика и астрономия}\newline \sIIbM & Механика, кинематика & Роботика, сензори за движение & mBot2, LEGO SPIKE Prime (Л2); безжични логери (Л3) & Позиция, скорост, ускорение & Регресионен модел в Jupyter (DH) & „Математика на движението“: проверка на законите на Нютон, модел – измерване \\
6 & \textbf{Информатика}\newline \sIIbIT\,\sIIbF\,\sIIbB & Програмиране, вградени системи & Микроконтролери, edge AI & ESP32/Arduino, Raspberry Pi~5 + AI HAT+ (Л2 → оранжерията) & Сензорни данни по MQTT & AI-2: ML модел на устройството (edge AI) & Автоматично поливане и вентилация в безопасни граници \\
7 & \textbf{Химия и ООС}\newline \sIIbI\,\sIIbF & Лабораторна автоматизация & Роботика, машинно зрение & Настолна роботизирана ръка с камера (Л2 ↔ Л3) & Цвят на разтвора, обем на титранта & AI-2 (лабораторни данни): CV за точката на еквивалентност & „Робот-лаборант“: автоматично титруване и сортиране на проби \\
8 & \textbf{География и икономика}\newline \sIIbM\,\sIIbI\,\sIIbB & Гео\-про\-стран\-ствен анализ & Дронове, фотограметрия & Дрон с камера за картиране, програмируеми дронове (Л2); AI сървър (DH) & Аеро\-фото\-снимки, GPS координати & Учебен CV модел за сегментация на покритието & „Дрон картира двора“: карта, площи, 3D модел, разходи за озеленяване \\
9 & \textbf{Информационни технологии}\newline \sIIbM\,\sIIbF\,\sIIbB\,\sIIbG & Дизайн и дигитално произ\-вод\-ство & 3D печат, лазерно рязане, 3D сканиране & 3D принтери, лазерен резач и 3D скенер (спец.) (Л2) & CAD модели, часове на печат, видео от камерата & AI-3: откриване на неуспешен печат & Корпуси за сензорите в оранжерията; релефни и биологични модели \\
\multicolumn{8}{@{}l}{\cellcolor{sand}\textbf{Направление „Математика и информатика“ – Л1 и Data Hub}}\\
10 & \textbf{Математика}\newline \sIIbG\,\sIIbF\,\sIIbI & Статистика и анализ на данни & Платформа за данни (ThingsBoard, времеви редове) & 26 лаптопа (Л1); NAS, Jupyter, Grafana (DH) & Исторически данни от оранжерията и метеостанцията & AI-2: предиктивен модел и оценка на грешката & „Климатът на училището“: статистически и предиктивен анализ \\
11 & \textbf{Информатика}\newline всички предмети & AI и машинно обучение & Локален LLM + RAG & AI сървър (RTX 5090, 32~GB) (DH); лаптопи (Л1) & Псевдонимизирани въпроси, код, протоколи; учебно съдържание & AI-1 „СМГ AI тютор“ & Помощ при програмиране, задачи по ниво; учителят одобрява (human-in-the-loop) \\
12 & \textbf{Информационни технологии}\newline \sIIbF\,\sIIbM & Енергийна ефективност, мрежи & IoT телеметрия, термография & UPS и мрежа (Л1), IoT „умна къща“ (Л2), ръчни термокамери (Л3), дисплеи (ВОСКС) & Консумация (kWh), часове работа, UPS логове, термограми & AI-3: прогноза на натоварването, аномалии & „Енергия в класната стая“: оптимизация, изключване на неизползвани дисплеи \\
\multicolumn{8}{@{}l}{\cellcolor{sand}\textbf{ВОСКС – 8 високотехнологични оборудвани и свързани класни стаи}}\\
13 & \textbf{География и икономика}\newline \sIIbF\,\sIIbM\,\sIIbIT & Науки за Земята, климат & IoT сензори & CO$_2$/температурни сензори в 8 ВОСКС; метеостанция (Л3) & CO$_2$, температура, влажност вътре и навън & AI-3: препоръка за проветряване & „Климатът на училището“: качество на въздуха спрямо външния климат \\
14 & \textbf{Всички 7 предмета} & Дигитално и хибридно обучение & Интерактивни дисплеи, видеовръзка & Дисплей с OPS, камера, документ-камера, учителски лаптоп (ВОСКС) & Табла на живо от Л1–Л3, видео от микроскопа и телескопа & AI-1 тютор: задачи по ниво & Урок с реални данни от лабораториите без преместване на класа \\
\midrule
\multicolumn{8}{@{}p{\textwidth}@{}}{\textbf{Покритие:} предмети в брой вериги – \sIIbM~11, \sIIbI~9, \sIIbIT~6, \sIIbF~11, \sIIbH~6, \sIIbB~8, \sIIbG~6; AI-1 – 3 вериги, AI-2 – 6, AI-3 – 3, учебни ML модели – 3; по направления – природни науки 4, роботика 5, математика и информатика 3, ВОСКС 2; обхванати са и осемте STEM сценария за първата година.}\\
\end{xltabular}
}

\subsection*{Приложение Б. Графична схема на връзките}
\addcontentsline{toc}{subsection}{Приложение Б. Графична схема на връзките}

\noindent\begin{minipage}{\textwidth}
\setlength{\parskip}{5pt}
Схемата свързва учебните предмети (цветни значки), технологиите (бели карета) и STEM пространствата (групи). В центъра е „СМГ STEM Data Hub“: всички пространства обменят данни само през него, а пунктираните връзки показват споделено оборудване между пространствата по общия график за резервации.

\begin{center}
\begin{tikzpicture}[
  font=\scriptsize,
  tech/.style={draw=grayline, fill=white, rounded corners=2pt, align=left, text width=3.45cm, minimum height=1.05cm, inner sep=2.5pt, execute at begin node={\hyphenpenalty=10000\exhyphenpenalty=10000}},
  ttl/.style={align=center, text width=7.4cm, inner sep=1pt},
  zone/.style={draw=grayline, rounded corners=4pt, inner sep=5pt},
  ai/.style={draw=navy, fill=white, rounded corners=2pt, align=left, text width=2.62cm, minimum height=1.45cm, inner sep=2.5pt, anchor=north, execute at begin node={\hyphenpenalty=10000\exhyphenpenalty=10000}},
  flow/.style={{Latex[length=2mm]}-{Latex[length=2mm]}, very thick, teal},
  share/.style={{Latex[length=1.8mm]}-{Latex[length=1.8mm]}, thick, dashed, navy},
  lbl/.style={font=\tiny, align=left, text=black!75}
]
% ---------- Горен ред: Л2 и Л3 ----------
\node[ttl] (t2) at (-4.1,7.05) {\textbf{\small Л2 „Роботика и мейкърспейс“}\\[-1pt]{\tiny Роботика и кибер-физични системи · 26 места + мейкър зона}};
\node[tech] (r1) at (-6.0,6.05) {\textbf{mBot2, LEGO SPIKE Prime, ROS2 робот}\\[2pt]\sIIbF\,\sIIbM\,\sIIbI};
\node[tech] (r2) at (-2.2,6.05) {\textbf{Роботизирана ръка с камера}\\[2pt]\sIIbH\,\sIIbI\,\sIIbF};
\node[tech] (r3) at (-6.0,4.85) {\textbf{Дронове: полет и картиране}\\[2pt]\sIIbG\,\sIIbM\,\sIIbI\,\sIIbB};
\node[tech] (r4) at (-2.2,4.85) {\textbf{ESP32/Arduino, Raspberry Pi~5 + AI HAT+}\\[2pt]\sIIbI\,\sIIbIT\,\sIIbF\,\sIIbB};
\node[tech] (r5) at (-6.0,3.65) {\textbf{3D принтери, лазерен резач, 3D скенер}\\[2pt]\sIIbIT\,\sIIbM\,\sIIbF\,\sIIbB\,\sIIbG};
\node[tech] (r6) at (-2.2,3.65) {\textbf{IoT комплект „умна къща“}\\[2pt]\sIIbIT\,\sIIbF\,\sIIbM};

\node[ttl] (t3) at (4.1,7.05) {\textbf{\small Л3 „Природни науки“ + умна оранжерия}\\[-1pt]{\tiny Природни науки · 24 места на 6 острова}};
\node[tech] (n1) at (2.2,6.05) {\textbf{Умна оранжерия: IoT сензори и камера}\\[2pt]\sIIbB\,\sIIbH\,\sIIbF\,\sIIbM\,\sIIbI};
\node[tech] (n2) at (6.0,6.05) {\textbf{Метеостанция}\\[2pt]\sIIbG\,\sIIbF\,\sIIbM\,\sIIbI\,\sIIbIT};
\node[tech] (n3) at (2.2,4.85) {\textbf{Логери, pH/EC метри, спектрофотометър}\\[2pt]\sIIbH\,\sIIbB\,\sIIbF\,\sIIbM};
\node[tech] (n4) at (6.0,4.85) {\textbf{Цифрови микроскопи, тринокулярен с камера}\\[2pt]\sIIbB\,\sIIbI\,\sIIbH};
\node[tech] (n5) at (2.2,3.65) {\textbf{GoTo телескоп с астрокамера}\\[2pt]\sIIbF\,\sIIbM\,\sIIbI};
\node[tech] (n6) at (6.0,3.65) {\textbf{Ръчни термокамери (Л3)}\\[2pt]\sIIbIT\,\sIIbF\,\sIIbM};

% ---------- Център: Data Hub ----------
\node[align=center, text=navy] (ht) at (0,1.85) {\textbf{\normalsize СМГ STEM Data Hub}};
\node[lbl, align=center] (hp) at (0,1.45) {MQTT → ThingsBoard → NAS → AI сървър → Grafana \ · \ Moodle \ · \ резервации и инвентар};
\node[ai] (h0) at (-4.41,1.18) {\textbf{AI сървър + NAS}\\{\tiny RTX 5090, 32 GB VRAM}\\[1pt]\sIIball};
\node[ai] (h1) at (-1.47,1.18) {\textbf{AI-1: AI тютор}\\{\tiny „СМГ AI тютор“: локален LLM + RAG, учителят одобрява}};
\node[ai] (h2) at (1.47,1.18) {\textbf{AI-2: анализ}\\{\tiny оранжерия и лабораторни данни: прогноза, аномалии, CV}};
\node[ai] (h3) at (4.41,1.18) {\textbf{AI-3: управление}\\{\tiny интелигентен график и прогнозна поддръжка}};

% ---------- Долен ред: Л1 и ВОСКС ----------
\node[ttl] (t1) at (-4.1,-1.8) {\textbf{\small Л1 „Данни и AI“}\\[-1pt]{\tiny Математика и информатика · 26 места}};
\node[tech] (m1) at (-6.0,-2.8) {\textbf{26 лаптопа: Jupyter, GeoGebra}\\[2pt]\sIIbM\,\sIIbI\,\sIIbF\,\sIIbG};
\node[tech] (m2) at (-2.2,-2.8) {\textbf{Мрежа, UPS, Raspberry Pi~5}\\[2pt]\sIIbIT\,\sIIbF\,\sIIbM};

\node[ttl] (t4) at (4.1,-1.8) {\textbf{\small 8 ВОСКС}\\[-1pt]{\tiny обикновени класни стаи в основната сграда}};
\node[tech] (v1) at (2.2,-2.8) {\textbf{Дисплей с OPS, камера, документ-камера}\\[2pt]\sIIball};
\node[tech] (v2) at (6.0,-2.8) {\textbf{CO$_2$ и температурни сензори}\\[2pt]\sIIbG\,\sIIbF\,\sIIbM\,\sIIbIT};

% ---------- Фонове на пространствата ----------
\begin{scope}[on background layer]
  \node[zone, fill=sand, fit=(t2)(r1)(r6)] (Z2) {};
  \node[zone, fill=sand, fit=(t3)(n1)(n6)] (Z3) {};
  \node[zone, fill=navy!7, draw=navy, very thick, fit=(ht)(hp)(h0)(h3)] (HUB) {};
  \node[zone, fill=sand, fit=(t1)(m1)(m2)] (Z1) {};
  \node[zone, fill=teal!8, fit=(t4)(v1)(v2)] (Z4) {};
\end{scope}

% ---------- Потоци към/от Data Hub ----------
\draw[flow] (Z2.south -| -4.1,0) -- (HUB.north -| -4.1,0);
\node[lbl, anchor=west] at ($(Z2.south -| -3.95,0)!0.5!(HUB.north -| -3.95,0)$) {↓ телеметрия, видео\\↑ ML модели за устройствата (edge)};
\draw[flow] (Z3.south -| 4.1,0) -- (HUB.north -| 4.1,0);
\node[lbl, anchor=west] at ($(Z3.south -| 4.25,0)!0.5!(HUB.north -| 4.25,0)$) {↓ сензорни данни, снимки\\↑ прогнози, аномалии};
\draw[flow] (Z1.north -| -4.1,0) -- (HUB.south -| -4.1,0);
\node[lbl, anchor=west] at ($(Z1.north -| -3.95,0)!0.5!(HUB.south -| -3.95,0)$) {↑ анализи, обучени модели\\↓ исторически данни, AI-1};
\draw[flow] (Z4.north -| 4.1,0) -- (HUB.south -| 4.1,0);
\node[lbl, anchor=west] at ($(Z4.north -| 4.25,0)!0.5!(HUB.south -| 4.25,0)$) {↑ CO$_2$, температура\\↓ табла на живо, AI-1};

% ---------- Споделено оборудване между пространствата ----------
\draw[share] (Z2.north -| -1.4,0) to[out=60,in=120] node[lbl, above, align=center, text width=7cm]{„Робот-лаборант“; ESP32 и 3D корпуси за сензорите в оранжерията} (Z3.north -| 1.4,0);
\draw[share] (Z1.south -| -1.4,0) to[out=-60,in=-120] node[lbl, below, align=center, text width=10cm]{мобилни комплекти (логери, роботи, дронове) и лаптопи по общия график за резервации} (Z4.south -| 1.4,0);
\end{tikzpicture}
\end{center}

{\scriptsize
\textbf{Легенда.} \sIIbM~математика \enspace \sIIbI~информатика \enspace \sIIbIT~информационни технологии \enspace \sIIbF~физика и астрономия\\
\sIIbH~химия и ООС \enspace \sIIbB~биология и ЗО \enspace \sIIbG~география и икономика \enspace рамка – STEM пространство (Л1–Л3, ВОСКС)\\
\tikz[baseline=-0.6ex]\draw[teal,very thick](0,0)--(0.7,0); поток на данни през Data Hub (MQTT/ThingsBoard) \qquad \tikz[baseline=-0.6ex]\draw[navy,thick,dashed](0,0)--(0.7,0); споделено оборудване между пространствата
\par}
\end{minipage}
\clearpage
% --- sections/a_s3.tex
\subsection*{Приложение В. Схеми на функциониране на AI приложенията}
\addcontentsline{toc}{subsection}{Приложение В. Схеми на функциониране на AI приложенията}

Всяка схема следва веригата от т.~3: Input → Data Processing → AI Model/\allowbreak Service → Output → Decision/\allowbreak Action. Тъмносините етапи са автоматичната обработка, а тюркоазеният – решението или действието, което е под човешки контрол или в зададени от учителя граници. Пунктирната стрелка е обратната връзка (Human-in-the-loop): одобренията и корекциите се връщат като нови входни данни. Долната лента посочва хардуера и софтуера от бюджета и учебните предмети, които ползват приложението.

% #1–#5: съдържание на петте етапа; #6: текст на обратната връзка; #7: инфраструктура и предмети; #8: височина на блоковете
\newcommand{\AppVPipeline}[8]{%
\begin{tikzpicture}[font=\scriptsize,
  hdr/.style={draw=navy, fill=navy, text=white, font=\scriptsize\bfseries, rounded corners=2pt,
              text width=2.9cm, inner sep=1.5pt, minimum height=0.55cm, align=center},
  cnt/.style={draw=navy, fill=sand, rounded corners=2pt, inner sep=1.5pt, anchor=north},
  arr/.style={-{Stealth[length=1.8mm]}, thick, navy},
  fb/.style={-{Stealth[length=1.8mm]}, thick, dashed, teal},
  infra/.style={draw=grayline, fill=white, rounded corners=2pt, text width=16.4cm, inner sep=2.5pt, align=left, anchor=north west}]
\node[hdr] (h1) at (1.503,0)  {Input};
\node[hdr] (h2) at (4.899,0)  {Data Processing};
\node[hdr] (h3) at (8.295,0)  {AI Model/Service};
\node[hdr] (h4) at (11.691,0) {Output};
\node[hdr, fill=teal, draw=teal] (h5) at (15.087,0) {Decision/Action};
\node[cnt] (c1) at ($(h1.south)+(0,-0.08)$) {\parbox[c][#8][t]{2.9cm}{\fontsize{7.5}{8.8}\selectfont\raggedright #1}};
\node[cnt] (c2) at ($(h2.south)+(0,-0.08)$) {\parbox[c][#8][t]{2.9cm}{\fontsize{7.5}{8.8}\selectfont\raggedright #2}};
\node[cnt] (c3) at ($(h3.south)+(0,-0.08)$) {\parbox[c][#8][t]{2.9cm}{\fontsize{7.5}{8.8}\selectfont\raggedright #3}};
\node[cnt] (c4) at ($(h4.south)+(0,-0.08)$) {\parbox[c][#8][t]{2.9cm}{\fontsize{7.5}{8.8}\selectfont\raggedright #4}};
\node[cnt, draw=teal, thick] (c5) at ($(h5.south)+(0,-0.08)$) {\parbox[c][#8][t]{2.9cm}{\fontsize{7.5}{8.8}\selectfont\raggedright #5}};
\draw[arr] (h1.east) -- (h2.west);
\draw[arr] (h2.east) -- (h3.west);
\draw[arr] (h3.east) -- (h4.west);
\draw[arr] (h4.east) -- (h5.west);
\node[fit=(c1)(c2)(c3)(c4)(c5), inner sep=0pt] (all) {};
\draw[fb] (c5.south) -- (c5.south |- all.south) -- ++(0,-0.3) -| node[pos=0.25, below=1pt, text=teal, font=\scriptsize\itshape] {Human-in-the-loop: #6} (c1.south);
\node[infra] at ($(all.south west)+(0,-0.85)$) {#7};
\end{tikzpicture}}

\medskip
\noindent\begin{minipage}{\textwidth}
\textbf{\color{navy}В.1. AI-1 „СМГ AI тютор“ – AI за обучение}\par\smallskip
\AppVPipeline
{\textbullet~Въпроси и код на учениците\\
 \textbullet~Лабораторни протоколи\\
 \textbullet~Учебни програми, ръководства и задачи на СМГ\\
 \textbullet~Ниво на ученика (под псевдоним)}
{\textbullet~Вход през Moodle с роля\\
 \textbullet~Имената се заменят с псевдоними\\
 \textbullet~Съдържанието се индексира във векторна база\\
 \textbullet~Логове в NAS до края на учебната година}
{\textbullet~Отворен LLM (Qwen/Llama) чрез Ollama/vLLM на AI станцията\\
 \textbullet~RAG + цитат към източника\\
 \textbullet~Правило „подсказка, не решение“}
{\textbullet~Подсказка стъпка по стъпка\\
 \textbullet~Задачи по ниво: базово, напреднало, олимпийско\\
 \textbullet~Чернова на обратна връзка по рубрика}
{\textbullet~Учителят одобрява задачите и оценката в Moodle\\
 \textbullet~Ученикът поправя кода или протокола\\
 \textbullet~Типичните грешки влизат в следващия урок}
{учителят одобрява или коригира; корекциите допълват базата от знания}
{\textbf{Работи върху:} AI работна станция (RTX 5090, 32~GB VRAM), NAS, Moodle, система за управление на класната стая. \textbf{Предмети:} информатика, информационни технологии, математика, физика и астрономия, химия, биология.}
{2.75cm}
\end{minipage}

\bigskip
\noindent\begin{minipage}{\textwidth}
\textbf{\color{navy}В.2. AI-2 „Умна оранжерия и лабораторни данни“ – AI за анализ на данни}\par\smallskip
\AppVPipeline
{\textbullet~Влажност на почвата, температура, светлина\\
 \textbullet~pH и EC на хранителния разтвор\\
 \textbullet~Камера над растенията\\
 \textbullet~Метеостанция и безжични логери}
{\textbullet~ESP32 → MQTT → ThingsBoard\\
 \textbullet~Почистване на липсващи стойности и шум\\
 \textbullet~5-минутни средни стойности\\
 \textbullet~Времеви редове и снимки в NAS}
{\textbullet~Prophet/LSTM: прогноза за 24 часа\\
 \textbullet~Isolation Forest: аномалии\\
 \textbullet~YOLO/MobileNet на Raspberry Pi AI HAT+ / Jetson: състояние на растенията}
{\textbullet~Прогноза за влажност и температура\\
 \textbullet~Сигнал за аномалия (напр. повреден сензор)\\
 \textbullet~Препоръка за поливане\\
 \textbullet~Табло в Grafana}
{\textbullet~Автоматично поливане в граници min/max\\
 \textbullet~Торове и химикали – само след потвърждение от учител\\
 \textbullet~Данни за уроци по 5 предмета}
{учителят задава границите и потвърждава дозирането; ръчен прекъсвач на помпата}
{\textbf{Работи върху:} сензори и умна оранжерия, метеостанция, безжични логери, Raspberry Pi AI HAT+, Jetson Orin Nano, ThingsBoard PE, NAS, Grafana, AI работна станция. \textbf{Предмети:} биология, химия, физика, информатика, математика (сценарии „Умна оранжерия“ и „AI под микроскоп“).}
{2.75cm}
\end{minipage}

\bigskip
\noindent\begin{minipage}{\textwidth}
\textbf{\color{navy}В.3. AI-3 „Интелигентен график и прогнозна поддръжка“ – AI за управление на средата}\par\smallskip
\AppVPipeline
{\textbullet~Резервации, разписание\\
 \textbullet~Часове печат на 3D принтерите\\
 \textbullet~Цикли на батериите на дроновете\\
 \textbullet~Логове на UPS, консумация\\
 \textbullet~Камери на 3D принтерите}
{\textbullet~Обединяване в Data Hub: ThingsBoard + резервации и инвентар\\
 \textbullet~Нормализиране по устройство и час\\
 \textbullet~Маркиране на инциденти}
{\textbullet~Прогноза на търсенето\\
 \textbullet~Модел за износване (Isolation Forest)\\
 \textbullet~Компютърно зрение за неуспешен печат\\
 \textbullet~Оптимизация на графика}
{\textbullet~Седмичен график: 3 лаборатории и 8 ВОСКС\\
 \textbullet~Сигнал за поддръжка\\
 \textbullet~Неизползвани дисплеи за изключване}
{\textbullet~Администраторът одобрява графика\\
 \textbullet~Печатът спира автоматично при дефект\\
 \textbullet~Поддръжка преди повредата}
{администраторът одобрява графика и поддръжката; учителят може да рестартира спрян печат}
{\textbf{Работи върху:} AI работна станция, ThingsBoard PE, система за резервации и инвентар, UPS, камери на 3D принтерите, Grafana. \textbf{Ползватели:} трите направления и 8-те ВОСКС (сценарий „Енергия в класната стая“).}
{2.95cm}
\end{minipage}
\clearpage
% --- sections/a_s5.tex
\subsection*{Приложение Г. RBS диаграма}
\addcontentsline{toc}{subsection}{Приложение Г. RBS диаграма}

Йерархия на рисковете на проекта (към раздел 5.1). Ниво 1 – девет категории по заданието; ниво 2 – конкретни рискове с код от регистъра (Приложение~Д). Цветът на всеки риск съответства на нивото му в матрицата: \colorbox{riskH}{\footnotesize висок (15–25)}, \colorbox{riskM}{\footnotesize среден (8–12)}, \colorbox{riskL}{\footnotesize нисък (1–6)}.

\begin{center}
\begin{forest}
for tree={
  grow'=0,
  draw=grayline, rounded corners=2pt,
  anchor=west, child anchor=west, parent anchor=east,
  l sep=7mm, s sep=1.2mm,
  inner sep=2.5pt,
  font=\footnotesize,
  edge={draw=navy!60, line width=0.6pt},
  edge path={\noexpand\path[\forestoption{edge}] (!u.parent anchor) -- +(3.5mm,0) |- (.child anchor)\forestoption{edge label};},
},
where level=0{fill=navy, text=white, text width=2.0cm, node options={align=center}, font=\footnotesize\bfseries, draw=navy}{},
where level=1{fill=sand, text width=2.75cm, node options={align=left}, font=\footnotesize\bfseries}{},
where level=2{text width=8.9cm, node options={align=left}}{},
[{Проект „STEM център в СМГ“}
  [Технологични (ТЕХ)
    [{ТЕХ-1 Несъвместимост на оборудване и софтуер от различни доставчици}, fill=riskM]
    [{ТЕХ-2 Повреда на ключово оборудване (AI станция, 3D принтери, роботи)}, fill=riskM]
    [{ТЕХ-3 Зависимост от доставчик (vendor lock-in) и остаряване}, fill=riskM]
  ]
  [Финансови (ФИН)
    [{ФИН-1 Поскъпване между концепцията и офертите по ЗОП}, fill=riskH]
    [{ФИН-2 Недопустими разходи и финансова корекция}, fill=riskM]
  ]
  [Организационни (ОРГ)
    [{ОРГ-1 Прекратена обособена позиция (без оферти, жалба в КЗК)}, fill=riskH]
    [{ОРГ-2 Забавени доставки – риск за срока 15.09.2027}, fill=riskH]
    [{ОРГ-3 Слаба координация между направленията}, fill=riskM]
  ]
  [Образователни (ОБР)
    [{ОБР-1 Неготовност и съпротива на учителите}, fill=riskH]
    [{ОБР-2 Ниско използване на оборудването в часовете}, fill=riskM]
  ]
  [AI
    [{AI-1 Неточни отговори и халюцинации на AI тютора}, fill=riskM]
    [{AI-2 Пристрастност (bias) в задачите и черновите на оценки}, fill=riskL]
    [{AI-3 Прекомерна зависимост на учениците от AI}, fill=riskM]
  ]
  [Данни (ДАН)
    [{ДАН-1 Загуба на данни в NAS}, fill=riskM]
    [{ДАН-2 Нарушена цялост на сензорните данни}, fill=riskM]
    [{ДАН-3 Изтичане на лични данни на непълнолетни}, fill=riskH]
  ]
  [Киберсигурност (КИБ)
    [{КИБ-1 Ransomware атака срещу AI сървъра и NAS}, fill=riskM]
    [{КИБ-2 Неоторизиран достъп до MQTT брокера и помпата}, fill=riskM]
    [{КИБ-3 Prompt injection в AI тютора}, fill=riskM]
  ]
  [Инфраструктурни (ИНФ)
    [{ИНФ-1 Недостиг на контакти и мощност на електроинсталацията}, fill=riskM]
    [{ИНФ-2 Недостатъчно покритие и капацитет на Wi-Fi}, fill=riskM]
    [{ИНФ-3 Влага и вентилация в сутерена (лазер, запояване)}, fill=riskL]
  ]
  [Интеграционни (ИНТ)
    [{ИНТ-1 Свързване на сензори от различни производители с Data Hub}, fill=riskH]
    [{ИНТ-2 Несъвместимост на Moodle, Microsoft 365, резервации и Data Hub}, fill=riskM]
    [{ИНТ-3 Слаба интердисциплинарност на сценариите}, fill=riskM]
  ]
]
\end{forest}
\end{center}

\subsection*{Приложение Д. Пълен регистър на рисковете (Risk Matrix)}
\addcontentsline{toc}{subsection}{Приложение Д. Пълен регистър на рисковете (Risk Matrix)}

Скали: \textbf{P} – 1 рядко, 2 малко вероятно, 3 възможно, 4 вероятно, 5 почти сигурно; \textbf{I} – 1 незначително, 2 малко, 3 умерено (забавяне до 2 седмици или до 1\,\% от бюджета), 4 значително (забавяне до 2 месеца, загуба на позиция или пряк ефект върху учениците), 5 критично (провал на цел, финансова корекция, нарушение на GDPR). \textbf{S = P × I}; нива: 1–6 нисък, 8–12 среден, 15–25 висок. Разпределение: 6 високи, 17 средни, 2 ниски.

\begin{center}
\begin{tikzpicture}[x=2.7cm,y=1.4cm]
\foreach \p in {1,...,5}{
  \foreach \i in {1,...,5}{
    \pgfmathtruncatemacro{\sc}{\p*\i}
    \ifnum\sc<7 \def\rc{riskL}\else\ifnum\sc<13 \def\rc{riskM}\else\def\rc{riskH}\fi\fi
    \fill[\rc] (\i-1,\p-1) rectangle (\i,\p);
    \draw[white, line width=1.5pt] (\i-1,\p-1) rectangle (\i,\p);
    \node[anchor=north east, font=\tiny, text=black!45, inner sep=1.5pt] at (\i,\p) {\sc};
  }
  \node[anchor=east, font=\footnotesize] at (0,\p-0.5) {\p};
}
\foreach \i in {1,...,5}{\node[anchor=north, font=\footnotesize] at (\i-0.5,0) {\i};}
\node[font=\footnotesize\bfseries] at (2.5,-0.55) {Въздействие (I)};
\node[font=\footnotesize\bfseries, rotate=90] at (-0.3,2.5) {Вероятност (P)};
\tikzset{rk/.style={font=\scriptsize, align=center, text width=2.6cm, yshift=-2.5pt}}
\node[rk] at (3.5,3.5) {ФИН-1, ОРГ-2\\ОБР-1, ИНТ-1};
\node[rk] at (2.5,3.5) {AI-1, AI-3\\ИНФ-2};
\node[rk] at (4.5,2.5) {ОРГ-1, ДАН-3};
\node[rk] at (3.5,2.5) {ОБР-2, КИБ-2\\ИНФ-1, ИНТ-3};
\node[rk, inner sep=0pt, yshift=1pt] at (2.5,2.5) {ТЕХ-1, ТЕХ-2\\ТЕХ-3, ОРГ-3\\ДАН-2, КИБ-3\\ИНТ-2};
\node[rk] at (4.5,1.5) {ФИН-2, КИБ-1};
\node[rk] at (3.5,1.5) {ДАН-1};
\node[rk] at (2.5,1.5) {AI-2, ИНФ-3};
\end{tikzpicture}
\end{center}

Отговорници (роли от раздел 6): РП – ръководител на проекта (зам.-директор УД); STEM – STEM Integration Specialist; AI/ML – AI/ML Specialist; Соф. – Software Developer; IoT – IoT/\allowbreak Robotics Engineer; Инфр. – Infrastructure Specialist (зам.-директор АСД и системен администратор); Киб. – Cybersecurity Specialist; ЗОП – юрист по ЗОП; Сч. – счетоводител.

\begingroup\footnotesize
\setlength{\tabcolsep}{3pt}
\renewcommand{\arraystretch}{1.05}
\setlength{\LTpre}{3pt}\setlength{\LTpost}{0pt}
\begin{xltabular}{\textwidth}{@{}>{\bfseries}l>{\raggedright\arraybackslash\hsize=0.85\hsize}X ccc >{\centering\arraybackslash}p{1.12cm}>{\raggedright\arraybackslash}p{1.05cm}>{\raggedright\arraybackslash\hsize=1.15\hsize}X@{}}
\toprule
\rowcolor{sand}Код & \textbf{Риск} & \textbf{P} & \textbf{I} & \textbf{S} & \textbf{Ниво} & \textbf{Отгов.} & \textbf{Стратегия: превантивни мерки}\\
\midrule
\endfirsthead
\toprule
\rowcolor{sand}Код & \textbf{Риск} & \textbf{P} & \textbf{I} & \textbf{S} & \textbf{Ниво} & \textbf{Отгов.} & \textbf{Стратегия: превантивни мерки}\\
\midrule
\endhead
\bottomrule
\endlastfoot
\rowcolor{navy!8}\multicolumn{8}{@{}l}{\textbf{Технологични рискове}}\\*
ТЕХ-1 & Несъвместимост на оборудване и софтуер от различни доставчици (драйвери, ОС, протоколи) & 3 & 3 & 9 & \cellcolor{riskM}Среден & Инфр. & \textit{Намаляване:} изисквания за съвместимост и отворени формати в техническите спецификации; приемни тестове при доставка\\
ТЕХ-2 & Повреда на ключово оборудване (AI станция, 3D принтери, роботи, логери) & 3 & 3 & 9 & \cellcolor{riskM}Среден & Инфр. & \textit{Прехвърляне:} гаранция ≥\,24 месеца със сервиз в България и срок за подмяна в договора; история на повредите в системата за инвентар\\
ТЕХ-3 & Зависимост от един доставчик (vendor lock-in) и технологично остаряване & 3 & 3 & 9 & \cellcolor{riskM}Среден & Соф. & \textit{Избягване:} отворени стандарти (MQTT, REST, CSV/JSON); софтуер с отворен код (Grafana, Node-RED, Ollama); спецификации „или еквивалент“\\
\rowcolor{navy!8}\multicolumn{8}{@{}l}{\textbf{Финансови рискове}}\\*
ФИН-1 & Поскъпване между концепцията (11.2026) и офертите по ЗОП (03–05.2027), особено на GPU и памети & 4 & 4 & 16 & \cellcolor{riskH}Висок & РП, Сч. & \textit{Намаляване:} пазарни консултации през 02.2027 и актуални прогнозни стойности; резерв в обхвата – подреден списък на позиции, чиито количества се намаляват\\
ФИН-2 & Недопустими разходи и финансова корекция (нарушение на ЗОП, надвишени лимити 45/20/10\,\%) & 2 & 5 & 10 & \cellcolor{riskM}Среден & ЗОП & \textit{Избягване:} двоен контрол на документацията; проверка на лимитите при всяка промяна на бюджета; съгласуване с НСЦ\\
\rowcolor{navy!8}\multicolumn{8}{@{}l}{\textbf{Организационни рискове}}\\*
ОРГ-1 & Прекратена обособена позиция: няма оферти, цени над прогнозната стойност, жалба в КЗК & 3 & 5 & 15 & \cellcolor{riskH}Висок & ЗОП & \textit{Намаляване:} три обособени позиции; спецификации без марки („или еквивалент“); резерв от $\approx$\,3,5~седмици в графика (21.08–14.09.2027); при липса на оферти – договаряне по чл.~79, ал.~1, т.~1 ЗОП\\
ОРГ-2 & Забавени доставки $\to$ неготовност за въвеждане в експлоатация на 15.09.2027 & 4 & 4 & 16 & \cellcolor{riskH}Висок & РП & \textit{Прехвърляне:} неустойки за забава в договорите; срок за доставка ≤\,8~седмици; ИКТ и AI се доставят първи; резерв 21.08–14.09.2027 след външното тестване\\
ОРГ-3 & Слаба координация между екипа и трите направления; претоварване на учителите & 3 & 3 & 9 & \cellcolor{riskM}Среден & РП & \textit{Намаляване:} комуникационен план: седмична оперативна среща, месечен отчет, общ канал в Teams; ясни роли (RACI)\\
\rowcolor{navy!8}\multicolumn{8}{@{}l}{\textbf{Образователни рискове}}\\*
ОБР-1 & Неготовност и съпротива на учителите при работа с оборудването и AI & 4 & 4 & 16 & \cellcolor{riskH}Висок & STEM & \textit{Намаляване:} обучения от доставчиците като част от договорите; вътрешно обучение 06–09.2027; наставник във всяко направление\\
ОБР-2 & Ниско използване на оборудването; трудно вграждане в учебните програми & 3 & 4 & 12 & \cellcolor{riskM}Среден & STEM & \textit{Намаляване:} 8 интердисциплинарни сценария с готови уроци; общ график на лабораториите; KPI за използване всеки срок\\
\rowcolor{navy!8}\multicolumn{8}{@{}l}{\textbf{AI рискове}}\\*
AI-1 & Неточни отговори и халюцинации на AI тютора & 4 & 3 & 12 & \cellcolor{riskM}Среден & AI/ML & \textit{Намаляване:} RAG с цитати; 200 контролни въпроса; учителско одобрение; външно тестване (3\,300~€, 2–20.08.2027) и 4-седмичен пилот\\
AI-2 & Пристрастност (bias) при генериране на задачи и чернови на оценки & 2 & 3 & 6 & \cellcolor{riskL}Нисък & AI/ML & \textit{Приемане:} човешки надзор (human-in-the-loop); проверка на извадки всеки срок; оценяване по рубрика\\
AI-3 & Прекомерна зависимост на учениците от AI (преписване, загуба на умения) & 4 & 3 & 12 & \cellcolor{riskM}Среден & РП & \textit{Намаляване:} политика за AI; режим „подсказка, не решение“; устна обосновка на проектите; без AI на контролни\\
\rowcolor{navy!8}\multicolumn{8}{@{}l}{\textbf{Рискове, свързани с данните}}\\*
ДАН-1 & Загуба на данни (отказ на диск в NAS, случайно изтриване) & 2 & 4 & 8 & \cellcolor{riskM}Среден & Инфр. & \textit{Намаляване:} RAID; архиви по правилото 3-2-1 с офлайн копие; тестово възстановяване всеки срок\\
ДАН-2 & Нарушена цялост на сензорните данни (дрейф, пропуски, грешни времеви печати) & 3 & 3 & 9 & \cellcolor{riskM}Среден & IoT & \textit{Намаляване:} калибриране по график; валидиране на входа в ThingsBoard; автоматично откриване на аномалии\\
ДАН-3 & Изтичане на лични данни на непълнолетни (код, оценки, снимки), вкл. чрез външни AI услуги & 3 & 5 & 15 & \cellcolor{riskH}Висок & Киб. & \textit{Избягване:} локален AI – данните не напускат училището; без лични данни във външни AI услуги; псевдонимизация; DPIA; достъп по роли; обучение по GDPR\\
\rowcolor{navy!8}\multicolumn{8}{@{}l}{\textbf{Рискове за киберсигурността}}\\*
КИБ-1 & Ransomware атака срещу AI сървъра и NAS & 2 & 5 & 10 & \cellcolor{riskM}Среден & Киб. & \textit{Намаляване:} защита на крайните устройства; сегментиране на мрежата (VLAN); редовни обновления; офлайн архив; план за възстановяване\\
КИБ-2 & Неоторизиран достъп до MQTT брокера и командите към помпата и вентилатора & 3 & 4 & 12 & \cellcolor{riskM}Среден & IoT & \textit{Намаляване:} TLS и удостоверяване на всяко устройство; ACL по теми; отделна IoT VLAN; логове (FMEA, F7)\\
КИБ-3 & Prompt injection и заобикаляне на ролевите ограничения на AI тютора & 3 & 3 & 9 & \cellcolor{riskM}Среден & Киб. & \textit{Намаляване:} филтри на входа и изхода; права по роли; логове; тест за сигурност от външния изпълнител\\
\rowcolor{navy!8}\multicolumn{8}{@{}l}{\textbf{Инфраструктурни рискове}}\\*
ИНФ-1 & Недостиг на контакти и мощност на електроинсталацията (3D принтери, лазер, AI станция) & 3 & 4 & 12 & \cellcolor{riskM}Среден & Инфр. & \textit{Намаляване:} аудит на електроинсталацията при инициирането; UPS и разклонители със защита; разпределение на товара; ремонт – със средства на училището\\
ИНФ-2 & Недостатъчно покритие и капацитет на Wi-Fi (26+ устройства в стая, IoT, сутерен) & 4 & 3 & 12 & \cellcolor{riskM}Среден & Инфр. & \textit{Намаляване:} радиообследване; точка за достъп във всяка лаборатория и ВОСКС; кабелна връзка за стационарните устройства\\
ИНФ-3 & Влага и недостатъчна вентилация в сутерена (лазер, запояване) & 2 & 3 & 6 & \cellcolor{riskL}Нисък & Инфр. & \textit{Приемане:} лазер с филтър; аспирация при запояване; наблюдение на влажността със сензори\\
\rowcolor{navy!8}\multicolumn{8}{@{}l}{\textbf{Интеграционни рискове}}\\*
ИНТ-1 & Трудно свързване на сензори и логери от различни производители с Data Hub (ThingsBoard, MQTT) & 4 & 4 & 16 & \cellcolor{riskH}Висок & IoT & \textit{Намаляване:} MQTT/REST API като изискване в спецификациите; прототип през 02.2027; интеграционен тест през 08.2027\\
ИНТ-2 & Несъвместимост между Moodle, Microsoft 365, системата за резервации и Data Hub (профили, единен вход) & 3 & 3 & 9 & \cellcolor{riskM}Среден & Соф. & \textit{Намаляване:} единна идентичност чрез акаунтите на МОН (Microsoft 365); стандартни API; поетапно внедряване\\
ИНТ-3 & Слаба интердисциплинарност – сценариите остават в рамките на един предмет & 3 & 4 & 12 & \cellcolor{riskM}Среден & STEM & \textit{Намаляване:} съвместно планиране на учителите; 8 интердисциплинарни сценария; STEM Integration Score всеки срок\\
\end{xltabular}
\endgroup

\subsection*{Приложение Е. FMEA таблица}
\addcontentsline{toc}{subsection}{Приложение Е. FMEA таблица}

Компонент: IoT система на умната оранжерия (сензори $\to$ ESP32 и Raspberry Pi $\to$ MQTT $\to$ ThingsBoard $\to$ AI модел $\to$ помпа и вентилатор). Скали 1–10: \textbf{I} (Impact) – 1 незабележимо, 5 прекъсване на учебна дейност или загуба на растения, 9–10 опасност за хора или имущество; \textbf{P} (Probability) – 1 практически невъзможно, 10 почти сигурно; \textbf{D} (Detection) – 1 откриване веднага и автоматично, 10 не може да се открие преди последствията. \textbf{RPN = I × P × D}; критичен отказ: RPN ≥\,150 или I ≥\,9 (оцветени редове). RPN′ – стойност след прилагане на мерките.

\begingroup\footnotesize
\setlength{\tabcolsep}{3pt}
\renewcommand{\arraystretch}{1.1}
\setlength{\LTpre}{3pt}\setlength{\LTpost}{0pt}
\begin{xltabular}{\textwidth}{@{}>{\bfseries}l>{\raggedright\arraybackslash\hsize=0.95\hsize}X>{\raggedright\arraybackslash\hsize=0.8\hsize}X cccc>{\raggedright\arraybackslash\hsize=1.25\hsize}X>{\centering\arraybackslash}p{1.2cm}@{}}
\toprule
\rowcolor{sand}№ & \textbf{Елемент: режим на отказ (причина)} & \textbf{Последствия} & \textbf{I} & \textbf{P} & \textbf{D} & \textbf{RPN} & \textbf{Мерки: П – превантивни, К – коригиращи} & \textbf{RPN′ (I×P×D)}\\
\midrule
\endfirsthead
\toprule
\rowcolor{sand}№ & \textbf{Елемент: режим на отказ (причина)} & \textbf{Последствия} & \textbf{I} & \textbf{P} & \textbf{D} & \textbf{RPN} & \textbf{Мерки: П – превантивни, К – коригиращи} & \textbf{RPN′ (I×P×D)}\\
\midrule
\endhead
\bottomrule
\endlastfoot
\rowcolor{riskH}F1 & \textbf{Сензор за влажност на почвата:} дрейф и корозия (соли, влага) & Грешни стойности $\to$ преполиване или изсушаване; грешни данни за уроците & 6 & 6 & 6 & 216 & П: капацитивни сензори, месечно калибриране, втори контролен сензор; К: подмяна, маркиране на данните като невалидни & 54 (6×3×3)\\
F2 & \textbf{ESP32, Wi-Fi:} загуба на Wi-Fi връзка (смущения, слаб сигнал, рестарт на точката за достъп) & Липсват данни; командите не пристигат & 5 & 6 & 3 & 90 & П: точка за достъп в лабораторията, буфериране на данните, локални правила; К: рестарт от watchdog, аларма „няма данни > 10 мин“ & 30 (5×3×2)\\
\rowcolor{riskH}F3 & \textbf{Помпа и реле:} помпата остава включена (залепнало реле, софтуерна грешка) & Наводнение, повреда на оборудване и под, опасност от токов удар & 9 & 3 & 5 & 135 & П: хардуерен таймер, резервоар с ограничен обем, „изключено“ по подразбиране; К: датчик за теч изключва помпата и праща известие & 20 (5×2×2)\\
F4 & \textbf{Захранване:} прекъсване на захранването & Спира контролерът; губят се буферирани данни; неопределено състояние при рестарт & 6 & 4 & 2 & 48 & П: UPS, безопасно състояние при стартиране; К: автоматично възстановяване, синхронизация на часовника (NTP) & 24 (6×2×2)\\
\rowcolor{riskH}F5 & \textbf{AI модел (AI-2):} грешна прогноза или препоръка (нетипично време, дрейф на модела, малко данни) & Неправилно поливане или вентилация; грешни изводи в уроците & 6 & 5 & 7 & 210 & П: безопасни граници извън модела, сравнение с базово правило, преобучение всеки срок; К: превключване към правила, преглед от учител & 72 (6×3×4)\\
F6 & \textbf{NAS, база с времеви редове:} загуба или повреда на историческите данни (отказ на диск, повредена база) & Загуба на данни за проекти и за обучение на моделите & 7 & 3 & 4 & 84 & П: RAID, ежедневен архив по правилото 3-2-1, наблюдение на дисковете (S.M.A.R.T.); К: възстановяване от архив & 28 (7×2×2)\\
\rowcolor{riskH}F7 & \textbf{MQTT брокер:} неоторизиран достъп (слаби пароли, отворен порт, липса на TLS) & Фалшиви данни; злонамерени команди към помпата и вентилатора & 9 & 3 & 7 & 189 & П: TLS, отделни идентификационни данни за всяко устройство, ACL по теми, IoT VLAN; К: логове, аларма при нетипични команди, блокиране на устройството & 54 (9×2×3)\\
F8 & \textbf{Камера:} отказ, замърсен обектив или лошо осветление & Моделът за компютърно зрение пропуска болести и вредители & 4 & 5 & 4 & 80 & П: автоматична проверка на кадъра (яркост, рязкост), седмично почистване; К: ръчен оглед, подмяна & 24 (4×3×2)\\
F9 & \textbf{Вентилатор:} блокиране или износване & Прегряване, висока влажност, мухъл & 6 & 4 & 4 & 96 & П: контрол на тока и оборотите, аларма при температура над прага; К: ръчен режим, резервен вентилатор & 24 (6×2×2)\\
F10 & \textbf{Raspberry Pi (шлюз):} повреда на microSD картата или файловата система (интензивен запис) & Спира шлюзът и локалното компютърно зрение & 5 & 5 & 5 & 125 & П: SSD вместо microSD, файлова система само за четене, логове към NAS; К: резервен образ за бързо възстановяване & 20 (5×2×2)\\
\midrule
\multicolumn{6}{@{}l}{\textbf{Сумарен RPN (10 режима)}} & \textbf{1\,273} & \textbf{Намаление след мерките: −72,5\,\%} & \textbf{350}\\
\end{xltabular}
\endgroup

Приоритет на действията (по RPN и тежест): F1 (216) $\to$ F5 (210) $\to$ F7 (189) $\to$ F3 (135, I\,=\,9) – изпълняват се до интеграционния тест през август 2027; следват F10, F9, F2, F6, F8 и F4. Тежестта (I) намалява само при промяна в конструкцията – при F3 резервоарът с ограничен обем и датчикът за теч свеждат последствията до малък локален теч (I: 9 $\to$ 5); при останалите режими мерките намаляват вероятността и подобряват откриването.
\clearpage
% --- sections/a_s6.tex
\subsection*{Приложение Ж. Организационна структура на проектния екип}
\addcontentsline{toc}{subsection}{Приложение Ж. Организационна структура на проектния екип}

Схемата показва слабата матрична структура на проектния екип (т.~6): плътните връзки са линия на подчиненост по проекта, пунктирните – консултиране, съгласуване и отчет. Плътна рамка – служители и ученици на СМГ в проектния екип; пунктирна рамка – външни участници и заинтересовани страни извън екипа.

\begin{center}
\begin{tikzpicture}[
  font=\footnotesize,
  every node/.append style={execute at begin node={\hyphenpenalty=10000\exhyphenpenalty=10000}},
  node distance=0.22cm,
  box/.style={draw=navy, fill=sand, rounded corners=2pt, align=center, text width=#1, inner sep=3pt, minimum height=0.95cm},
  box/.default=4.8cm,
  head/.style={draw=navy, fill=navy, text=white, rounded corners=2pt, align=center, text width=#1, inner sep=3pt, minimum height=0.8cm},
  head/.default=4.8cm,
  ext/.style={draw=teal, dashed, thick, fill=white, rounded corners=2pt, align=center, text width=#1, inner sep=3pt, minimum height=0.95cm},
  ext/.default=4.8cm,
  stud/.style={draw=navy, fill=white, rounded corners=2pt, align=center, text width=#1, inner sep=3pt, minimum height=0.8cm},
  stud/.default=4.8cm,
  rel/.style={draw=navy, thick},
  con/.style={draw=teal, thick, dashed},
]
% --- управление ---
\node[head=4.9cm] (sp) at (0,0) {\textbf{Спонсор}\\ директор на СМГ (възложител по ЗОП)};
\node[ext=4.3cm] (ps) at (-5.7,0) {\textbf{Педагогически съвет}\\ одобрява концепцията};
\node[ext=4.3cm] (nsc) at (5.7,0) {\textbf{НСЦ / МОН}\\ съгласуване и отчет по НПВУ};
\node[head=4.9cm] (pm) at (0,-1.75) {\textbf{Ръководител на проекта}\\ зам.-директор УД (≈\,30\,\% от времето)};
\node[box=4.3cm] (zop) at (5.7,-1.75) {\textbf{ЗОП и финанси}\\ юрист по ЗОП, счетоводител};
\node[ext=4.3cm] (par) at (-5.7,-1.75) {\textbf{Ученици и родители}\\ обратна връзка, политика за AI};

\draw[rel] (sp) -- (pm);
\draw[con] (ps) -- (sp);
\draw[con] (nsc) -- (sp);
\draw[rel] (pm) -- (zop);
\draw[con] (par) -- (pm);

% --- три блока ---
\node[head] (g1) at (-5.6,-3.65) {\textbf{Образователен блок}};
\node[head] (g2) at (0,-3.65) {\textbf{AI и софтуер}};
\node[head] (g3) at (5.6,-3.65) {\textbf{Инфраструктура, IoT и роботика}};

\draw[rel] (pm.south) -- ++(0,-0.5) coordinate (bus);
\draw[rel] (bus -| g1.north) -- (bus -| g3.north);
\foreach \g in {g1,g2,g3} \draw[rel,-{Stealth[length=2mm]}] (bus -| \g.north) -- (\g.north);

% Образователен блок
\node[box, below=of g1] (stem) {\textbf{STEM Integration Specialist}\\ учител по физика};
\node[box, below=of stem] (tch) {\textbf{Teachers / Subject Experts}\\ математика, биология, химия, география};
\node[box, below=of tch] (da) {\textbf{Data Analyst}\\ учител по математика (KPI, STEM Integration Score)};

% AI и софтуер
\node[box, below=of g2] (ai) {\textbf{AI/ML Specialist}\\ учител по информатика};
\node[ext, below=of ai] (aic) {\textbf{Външен AI консултант}\\ модели, RAG, оценка на точността};
\node[box, below=of aic] (sw) {\textbf{Software Developer}\\ учител по ИТ („СМГ STEM Data Hub“)};
\node[stud, below=of sw] (st) {ученици от Клуб „Роботика“ (под надзор)};

% Инфраструктура, IoT и роботика
\node[box, below=of g3] (iot) {\textbf{IoT/Robotics Engineer}\\ ръководител на Клуб „Роботика“};
\node[box, below=of iot] (inf) {\textbf{Infrastructure Specialist}\\ зам.-директор АСД + системен администратор};
\node[ext, below=of inf] (sec) {\textbf{Cybersecurity Specialist}\\ външен експерт по договор};

% рамки на блоковете
\begin{scope}[on background layer]
  \node[draw=grayline, fill=sand!40, rounded corners=4pt, inner sep=4pt, fit=(g1)(stem)(tch)(da)] (b1) {};
  \node[draw=grayline, fill=sand!40, rounded corners=4pt, inner sep=4pt, fit=(g2)(ai)(aic)(sw)(st)] (b2) {};
  \node[draw=grayline, fill=sand!40, rounded corners=4pt, inner sep=4pt, fit=(g3)(iot)(inf)(sec)] (b3) {};
\end{scope}

% външно тестване на AI
\node[ext=15.6cm, minimum height=0.8cm, anchor=north] (test) at ([yshift=-0.55cm]b2.south) {\textbf{Независим външен изпълнител – приемно тестване на AI интеграцията:} 200 контролни въпроса, сигурност (prompt injection, достъп по роли), натоварване с 30 потребители, верига сензор → AI → действие; 10 човекодни, 3\,300~€, 2–20 август 2027~г. + 4-седмичен пилот (15.09–12.10.2027~г.)};
\foreach \b in {b1,b2,b3} \draw[con,-{Stealth[length=2mm]}] (\b.south) -- (\b.south |- test.north);
\end{tikzpicture}
\end{center}

{\footnotesize Стрелките към външния изпълнител означават: AI/ML специалистът и софтуерният разработчик предоставят средата за тест, IoT инженерът – веригата в оранжерията, учителите – пилота, а специалистът по киберсигурност и анализаторът на данни задават критериите за приемане (≥\,90\,\% верни отговори, 0 критични уязвимости, време за отговор < 5~s при 30 едновременни потребители).\par}

\medskip
\noindent\begin{minipage}{\textwidth}
\textbf{Матрица RACI по основните работни пакети от WBS}\par\smallskip
{\footnotesize
\setlength{\tabcolsep}{2pt}
\renewcommand{\arraystretch}{1.12}
\newcommand{\Ar}{\cellcolor{navy!18}\textbf{A}}
\newcommand{\Rr}{\cellcolor{teal!18}\textbf{R}}
\begin{tabularx}{\textwidth}{@{}>{\raggedright\arraybackslash}X*{11}{>{\centering\arraybackslash}p{0.88cm}}@{}}
\toprule
\textbf{Работен пакет (период)} & \scriptsize\textbf{Дир.} & \scriptsize\textbf{РП} & \scriptsize\textbf{STEM} & \scriptsize\textbf{AI/ML} & \scriptsize\textbf{Соф.} & \scriptsize\textbf{IoT} & \scriptsize\textbf{Инфр.} & \scriptsize\textbf{Киб.} & \scriptsize\textbf{Учит.} & \scriptsize\textbf{Данни} & \scriptsize\textbf{ЗОП}\\
\midrule
Концепция, решение на Педагогическия съвет, съгласуване с НСЦ (ноем.–дек. 2026) & \Ar & \Rr & C & C & I & C & C & I & C & I & I\\
Технически спецификации по направления (ян.–февр. 2027) & I & \Ar & \Rr & \Rr & C & \Rr & \Rr & C & \Rr & I & C\\
Обществена поръчка по ЗОП, 3 обособени позиции (март–май 2027) & \Ar & \Rr & C & C & I & C & C & C & C & I & \Rr\\
Доставка, монтаж, мрежа и инвентар (юни–юли 2027) & I & \Ar & I & I & C & \Rr & \Rr & C & I & I & C\\
Интеграция на Data Hub, IoT и AI (юли–август 2027) & I & \Ar & C & \Rr & \Rr & \Rr & C & C & I & C & I\\
Външно тестване на AI и пилот (август–октомври 2027) & I & \Ar & C & \Rr & C & C & I & C & \Rr & C & C\\
Обучение на учители и STEM сценарии (юни–септември 2027) & I & \Ar & \Rr & C & I & C & I & I & \Rr & C & I\\
Мониторинг, KPI и STEM Integration Score (септ. 2027 – юни 2028) & I & \Ar & C & C & C & I & I & I & C & \Rr & I\\
Закриване: финален отчет и поуки (юли 2028) & \Ar & \Rr & C & C & I & I & I & I & C & C & \Rr\\
\bottomrule
\end{tabularx}\par}
\smallskip
{\footnotesize R – изпълнява; A – носи крайната отговорност (точно един за пакет); C – консултира; I – информира се. Дир. – директор (спонсор); РП – ръководител на проекта; STEM – специалист по STEM интеграция; Соф. – софтуерен разработчик; Инфр. – инфраструктура; Киб. – киберсигурност; Учит. – учители по предмети; Данни – анализатор на данни; ЗОП – юрист по ЗОП и счетоводител.\par}
\end{minipage}
\clearpage
% --- sections/a_s7.tex
\subsection*{Приложение З. AI Project Log – регистър на използването на AI}
\addcontentsline{toc}{subsection}{Приложение З. AI Project Log}

Инструмент във всички записи: Claude Code (Anthropic, модел Claude Opus~5.5). При „агенти“ задачата е разделена между паралелни подагенти: изпълнител и независим проверяващ. Времената са от логовете на работните потоци (ноември 2026~г.).

{\scriptsize
\begin{xltabular}{\textwidth}{@{}>{\raggedright\arraybackslash}p{0.3cm} >{\raggedright\arraybackslash}p{2.2cm} >{\raggedright\arraybackslash}X >{\raggedright\arraybackslash}X >{\raggedright\arraybackslash}p{1.55cm} >{\raggedright\arraybackslash}X@{}}
\toprule
\textbf{№} & \textbf{Задача} & \textbf{Prompt (заявка, съкратено)} & \textbf{Генериран резултат} & \textbf{Приет / отхвърлен} & \textbf{Проверка и оценка} \\
\midrule
\endhead
1 & Анализ на заданието, насоките и образеца & „Прочети заданието, насоките на НПВУ и Excel образеца; изведи изискванията, лимитите и структурата на листовете.“ & Списък на 9-те точки и изходните резултати; редове на листовете; лимити; 7 полета за отметки в образеца & Приет & Сверено с PDF и с XML на образеца; открито, че библиотеката openpyxl губи полетата за отметки, затова е избрано запазване на ниво XML \\
2 & Уточняване на решенията & AI формулира 4 въпроса с варианти (направления, СМР, ДДС, лимити) & Решенията на студента: математика и информатика, роботика, природни науки; без СМР; цена ÷ 1,2 и офертата е нетна цена; лимитите от таблицата в € & Приет & Решенията са на студента, а не на AI (human-in-the-loop) \\
3 & Проучване на цените (106 позиции) & „За всяка позиция избери модел, наличен при два български търговеца; при липса – оферта от чужбина; зареди всяка страница; не измисляй цени.“ & 10 агента, 48~мин паралелно: модели, линкове, цени, спецификации на български & Приет с корекции & 10 независими проверяващи агента заредиха отново всеки линк: 87 потвърдени, 19 коригирани \\
4 & Допълнения за ВОСКС & „Безжична система за презентиране, микрофон и графичен таблет за 8 класни стаи.“ & 3 позиции, 24~мин & Приет с корекция & Проверяващият коригира 2 от 3 позиции: замени офертата от Германия с модел, продаван от двама български търговци, и махна непотвърдени характеристики на микрофона \\
5 & Независимост на източниците & „Намери втори търговец, който не е от същата мрежа (stem1.bg / office1.bg).“ & Нови източници за 4 комплекта по природни науки (Laborbio, Schoolmarket, Store.bg), 19~мин & Приет & Независим агент повторно зареди страниците: 4 от 4 потвърдени \\
6 & Правила за оферти & „Изчисли цена в чужбина + доставка + мито + ДДС при внос + такса; × 1,35.“ & Формули в колони O–AB за всеки ред & Приет & Формулите са проверени на ръка за ЕС, извън ЕС и цифров софтуер \\
7 & Оптимизиране на количествата & „Доближи сумите до лимитите, без да ги надвишаваш, като вдигаш количествата равномерно.“ & STEM \NumSTEM{} и ВОСКС \NumVOSKS & Приет & Формулите в Excel са пресметнати независимо и съвпадат с Python до цент; при ръчна проверка (J45 $\neq$ J43 + J44) е открита и поправена грешка в генерирания изчислител \\
8 & Съкращаване на софтуер и специализирано оборудване & AI предложи Labster, хостван Snipe-IT и компактен AI компютър (GX10) & 3 позиции, общо около 24\,000~€ & Отхвърлени & Софтуерът щеше да надхвърли 10\,\%, а специализираното оборудване – 27\,000~€; GX10 няма текущ български търговец (на страницата продуктът е „изчерпан“) \\
9 & Чернови на раздели 1–6 и 8 & „Напиши раздела по заданието и листа с факти; дължина 1–1,5 страници; препратки към предишните лекции.“ & 7 раздела в LaTeX + приложения А–Ж, 18,5~мин & Приет с корекции & 7 рецензиращи агента: съгласуваност между разделите, преизчислени таблици, проверка на дължината с компилация \\
10 & STEM Integration Score & Зададени тегла и оценки; „изчисли и направи анализ на чувствителността“ & Матрица 9 $\times$ 4 и 2 сценария на чувствителност & Приет & Всички произведения и суми са преизчислени ръчно \\
11 & Текстове 6.1–6.3 в образеца & „Опиши идеята, обосновката и устойчивостта, до 500 думи всяко.“ & 409 / 363 / 307 думи & Приет & Сверено с насоките на НПВУ и с раздели 1 и 9 \\
12 & Презентация и HTML преглед & „По един слайд за всяка точка, 10–12 минути; HTML преглед на Excel файла.“ & PPTX, PDF и HTML & Приет & Визуален преглед на всеки слайд и на HTML страницата \\
\bottomrule
\end{xltabular}}


\end{document}
